\chapter{Current Mirrors}\label{ch:Current_Mirrors}

This chapter covers various types of current mirrors,
which are circuits designed to copy (or ``mirror'') a current
from one its branch to another (amplifying or attenuating
it the process).

Ideally a current mirror should have perfect current
replication, infinite output resistance, zero input
resistance and no voltage drop across the mirror. However,
real current mirrors deviate from the ideal due to various
non-ideal effects, including component tolerances, component
parasitics, loading effects, channel length modulation,
base width modulation and temperature variations.

Channel Length Modulation and Base Width Modulation are
worth a discussion.
\begin{enumerate}
  \item \textbf{(MOSFET) Channel Length Modulation:} Generally
    we assume $\lambda = 0$ in MOSFET equations when analyzing
    transistor behaviour for simplicity. In real circuits,
    however, this is not the case. When $\lambda$ is non-zero,
    the drain current varies with changes in the drain-source
    voltage by a factor of $(1 + \lambda V_{DS})$. This affects
    the accuracy of current mirroring as the output current
    will vary with changes in the drain to source voltage.
    Cascoding may help reduce the impact of channel length modulation.
  \item \textbf{(BJT) Base Width Modulation:} Similar to channel length
    modulation in MOSFETs, BJTs also exhibit a phenomenon
    known as base width modulation (or Early effect). This
    effect causes the collector current to vary with changes
    in the collector-emitter voltage, leading to inaccuracies
    in current mirroring.
\end{enumerate}

All mirrors in this chapter are implemented using N-Channel
MOSFETs. It is possible to eadapt the configurations to
BJTs, or other transistor types regardless of technology.
Refer to~\ref{sec:theory_general_complementary} for more
information on converting circuits to their complementary
forms.

The primary use of current mirrors is to provide a stable
current source or sink in analog circuits. They can also
be used to steer currents into different branches of a
circuit, or to replace resistive loads in amplifiers and
other circuits thereby acting as active loads.

\section{Topologies}
\subsection{Widlar Mirror}

\subsubsection{Overview}
Constructed from two transistors, It is the
simplest form of a current mirror that can provide a
current gain different from unity.

\subsubsection{Schematic}
\begin{circuitfig}
  \draw (2.98, 7.905) to[american voltage source, l_={$V_{cc}$}] (1.73, 7.905);
	\node[nmos, arrowmos, xscale=-1] (N1) at (3.25, 4.405){} node[anchor=east] at (N1.text){$M_0$};
	\node[nmos, arrowmos] (N2) at (5.75, 4.405){} node[anchor=west] at (N2.text){$M_1$};
	\draw (3.25, 7.405) to[european resistor, l={$G_2$}] (3.25, 5.905);
	\draw (5.75, 5.655) -| (5.75, 5.175);
	\draw (4.5, 4.405) |- (3.25, 5.405);
	\draw (4.77, 4.405) -- (4.23, 4.405);
	\draw (3.25, 3.635) -| (3.25, 3.405) -| (5.75, 3.635);
	\draw (3.75, 2.655) to[american voltage source, l={$V_{ee}$}] (2.5, 2.655);
	\draw (3.75, 2.655) -| (4.5, 3.405);
	\node[ground, rotate=-90] at (1.73, 7.905){};
	\node[ground, rotate=-90] at (2.5, 2.655){};
	\draw[dash pattern={on 1.6pt off 0.8pt}, -stealth] (5.75, 5.655) -- (6.25, 5.655);
	\node[shape=rectangle, minimum width=0.965cm, minimum height=0.465cm] (N3) at (6.75, 5.615){} node[anchor=center] at (N3.text){$V_{out}$};
	\draw (5.75, 7.365) to[european resistor, l={\textcolor{rgb,255:red,255;green,0;blue,0}{$G_3$}}] (5.75, 5.865);
	\draw (5.75, 5.655) -- (5.75, 5.885);
	\draw (3.25, 7.405) |- (2.98, 7.905);
	\draw (3.25, 7.905) -| (5.75, 7.365);
	\node[shape=rectangle, minimum width=4.965cm, minimum height=1.09cm] at (4.375, 9.928){} node[anchor=north, align=center, text width=4.577cm, inner sep=6pt] at (4.375, 10.49){Large Signal Schematic\\(Sink)};
	\node[shape=rectangle, draw={rgb,255:red,255;green,0;blue,0}, line width=1pt, dash pattern={on 4pt off 2pt on 1pt off 2pt on 1pt off 2pt}, minimum width=1.465cm, minimum height=2.195cm] at (6, 7){};
	\node[shape=rectangle, minimum width=3.09cm, minimum height=2.465cm] at (8.438, 6.865){} node[anchor=north west, align=left, text width=2.702cm, inner sep=6pt] at (6.875, 8.115){\textcolor{rgb,255:red,255;green,0;blue,0}{\footnotesize Exists Only For Analysis. In  Implementation take output directly from V\_\{out\}}};
	\node[shape=rectangle, minimum width=0.965cm, minimum height=0.465cm] (N4) at (2.25, 5.405){} node[anchor=center] at (N4.text){$V_{x}$};
	\draw[dash pattern={on 1.6pt off 0.8pt}, -stealth] (3.25, 5.405) |- (2.75, 5.405);
	\node[shape=rectangle, minimum width=0.965cm, minimum height=0.465cm] (N5) at (4.5, 6.115){} node[anchor=center] at (N5.text){$I_{G}$};
	\draw (3.25, 5.175) -| (3.25, 5.405);
	\draw (3.25, 5.405) -| (3.25, 5.905);
	\draw (10, 1.75) |- (10, 11);
	\draw (17.25, 5.563) to[american current source, l_={$g_{m_1}v_{gs}$}] (17.25, 4.063);
	\draw (18.75, 5.563) to[european resistor, l_={$g_{o_1}$}] (18.75, 4.063);
	\draw (17.25, 4.063) -- (18.75, 4.063);
	\draw (18.75, 5.563) -- (17.25, 5.563);
	\draw (17.25, 4.063) -- (16, 4.063);
	\draw (16, 5.563) -- (15.25, 5.563);
	\node[shape=rectangle, minimum width=0.465cm, minimum height=0.465cm] at (16, 5.313){} node[anchor=center, align=center, text width=0.077cm, inner sep=6pt] at (16, 5.313){\footnotesize G};
	\node[shape=rectangle, minimum width=0.465cm, minimum height=0.465cm] at (16, 4.313){} node[anchor=center, align=center, text width=0.077cm, inner sep=6pt] at (16, 4.313){\footnotesize S};
	\node[shape=rectangle, minimum width=0.465cm, minimum height=0.465cm] at (19.25, 5.313){} node[anchor=center, align=center, text width=0.077cm, inner sep=6pt] at (19.25, 5.313){\footnotesize D};
	\draw (15.25, 5.563) -- (14.5, 5.563);
	\draw (13, 5.563) to[american current source, l={$g_{m_0}v_{gs}$}] (13, 4.063);
	\draw (11.5, 5.563) to[european resistor, l={$g_{o_0}$}] (11.5, 4.063);
	\draw (13, 5.563) -- (11.5, 5.563);
	\draw (13, 4.063) -- (11.5, 4.063);
	\draw (14.5, 4.063) -- (13, 4.063);
	\node[shape=rectangle, minimum width=0.465cm, minimum height=0.465cm] at (14.5, 5.313){} node[anchor=center, align=center, text width=0.077cm, inner sep=6pt] at (14.5, 5.313){\footnotesize G};
	\node[shape=rectangle, minimum width=0.465cm, minimum height=0.465cm] at (14.5, 4.313){} node[anchor=center, align=center, text width=0.077cm, inner sep=6pt] at (14.5, 4.313){\footnotesize S};
	\node[shape=rectangle, minimum width=0.465cm, minimum height=0.465cm] at (11, 5.313){} node[anchor=center, align=center, text width=0.077cm, inner sep=6pt] at (11, 5.313){\footnotesize D};
	\draw (11.5, 7.563) to[european resistor, l={$G_2$}] (11.5, 6.063);
	\draw (11.5, 6.063) -| (11.5, 5.562);
	\draw (11.5, 7.563) -| (11.5, 8.062) -- (12.25, 8.062) -| (12.25, 7.812);
	\draw (18.75, 7.563) to[european resistor, l={\textcolor{rgb,255:red,255;green,0;blue,0}{$G_3$}}] (18.75, 6.063);
	\draw (18.75, 6.063) -| (18.75, 5.562);
	\draw (18.75, 7.563) -| (18.75, 8.062) |- (18, 8.062) -| (18, 7.812);
	\draw[dash pattern={on 1.6pt off 0.8pt}, -stealth] (18.75, 5.563) -- (19.25, 5.563);
	\node[shape=rectangle, minimum width=0.715cm, minimum height=0.465cm] (N6) at (10.625, 5.563){} node[anchor=center] at (N6.text){$V_{x}$};
	\node[shape=rectangle, minimum width=0.965cm, minimum height=0.465cm] (N7) at (19.75, 5.563){} node[anchor=center] at (N7.text){$V_{out}$};
	\node[ground] at (14.5, 4.063){};
	\node[ground] at (16, 4.063){};
	\node[ground] at (12.25, 7.812){};
	\node[ground] at (18, 7.812){};
	\node[shape=rectangle, minimum width=0.965cm, minimum height=0.465cm] (N8) at (15.25, 6.813){} node[anchor=center] at (N8.text){$I_{G}$};
	\draw[dash pattern={on 1.6pt off 0.8pt}, -stealth] (11.5, 5.563) -- (11, 5.563);
	\node[shape=rectangle, minimum width=4.215cm, minimum height=1.09cm] at (15.375, 9.937){} node[anchor=north, align=center, text width=3.827cm, inner sep=6pt] at (15.375, 10.5){Small Signal Schematic\\(Sink)};
	\node[circ] at (3.25, 5.405){};
	\node[circ] at (4.5, 4.405){};
	\node[circ] at (4.5, 3.405){};
	\node[circ] at (3.25, 7.905){};
	\draw (13, 5.563) |- (15.25, 6.063) -| (15.25, 5.563);
	\draw[dash pattern={on 1.6pt off 0.8pt}, -stealth] (15.25, 6.063) -- (15.25, 6.563);
	\draw[dash pattern={on 1.6pt off 0.8pt}, -stealth] (4.5, 5.365) -- (4.5, 5.865);
	\node[circ] at (11.5, 5.563){};
	\node[circ] at (15.25, 5.563){};
	\node[circ] at (18.75, 5.563){};
	\node[circ] at (13, 5.563){};
	\node[shape=rectangle, draw={rgb,255:red,255;green,0;blue,0}, line width=1pt, dash pattern={on 4pt off 2pt on 1pt off 2pt on 1pt off 2pt}, minimum width=1.465cm, minimum height=1.965cm] at (19, 7.063){};
\end{circuitfig}

\subsubsection{Transfer Function}
Here's a table of the voltages at each terminal of the transistors:

\begin{tabularx}{\linewidth}{|X|X|X|X|X|X|} \hline
  Transistor & Gate (G) & Source (S) & Drain (D) & $V_{gs}$ & $V_{ds}$ \\ \hline
  $M_0$ & $V_{x}$ & $V_{ee}$ & $V_{x}$ & $V_{x} - V_{ee}$ & $V_{x} - V_{ee}$ \\ \hline
  $M_1$ & $V_{x}$ & $V_{ee}$ & $V_{out}$ & $V_{x} - V_{ee}$ & $V_{out} - V_{ee}$ \\ \hline
\end{tabularx}

\begin{enumerate}
  \item \textbf{Large Signal Analysis}:\\
    Applying KCL at $V_{out}$:
    \begin{align*}
      (V_{out} - V_{cc})G_3 + I_{M_1} &= 0 \\
      V_{out}G_3 - V_{cc}G_3 + I_{M_1} &= 0 \\
      V_{out} &= V_{cc} - \frac{I_{M_1}}{G_3} \tag{1}\end{align*}
    Applying KCL at $V_{x}$:
    \begin{align*}
      (V_{x} - V_{cc})G_2 + I_{M_0} + I_{G} &= 0 \\
      V_{x}G_2 - V_{cc}G_2 + (I_{M_0} + I_{G}) &= 0 \\
      V_{x}G_2 &= - (I_{M_0} + I_{G}) + V_{cc}G_2 \\
      V_{x} &= V_{cc} - \frac{(I_{M_0} + I_{G})}{G_2} \tag{2}
    \end{align*}
    Using the transistor equations for $I_{M_0}$ and
    $I_{M_1}$:
    \begin{align*}
      I_{M_0} &= k_{n_0}(V_{gs_0} - V_{th})^2 \cdot (1 + \lambda V_{ds_0}) \\
              &= k_{n_0}(V_{x} - V_{ee} - V_{th})^2 \cdot (1 + \lambda (V_{x} - V_{ee})) \\
              &= k_{n_0}\left(\left(V_{cc} - \frac{(I_{M_0} + I_{G})}{G_2}\right) - V_{ee} - V_{th}\right)^2 \cdot \left(1 + \lambda \left(\left(V_{cc} - \frac{(I_{M_0} + I_{G})}{G_2}\right) - V_{ee}\right)\right) \\
              &= k_{n_0}\left(\left(V_{cc} - \frac{(I_{M_0} + I_{G})}{G_2}\right) - V_{ee} - V_{th}\right)^2 \quad \text{[assuming } \lambda = 0 \text{]} \tag{3} \\ \\
      I_{M_1} &= k_{n_1}(V_{gs_1} - V_{th})^2 \cdot (1 + \lambda V_{ds_1}) \\
              &= k_{n_1}(V_{x} - V_{ee} - V_{th})^2 \cdot (1 + \lambda (V_{out} - V_{ee})) \\
              &= k_{n_1}\left(\left(V_{cc} - \frac{(I_{M_0} + I_{G})}{G_2}\right) - V_{ee} - V_{th}\right)^2 \cdot \left(1 + \lambda \left(V_{cc} - \frac{I_{M_1}}{G_3} - V_{ee}\right)\right) \\
              &= k_{n_1}\left(\left(V_{cc} - \frac{(I_{M_0} + I_{G})}{G_2}\right) - V_{ee} - V_{th}\right)^2 \quad \text{[assuming } \lambda = 0 \text{]} \tag{4} \\
    \end{align*}
    It can be seen that,
    \begin{align*}
      I_{in}  = I_{M_0} &\approx I_{M_1}  \quad \text{[if } k_{n_0} = k_{n_1} \text{]} \\
      I_{out} = I_{M_1}
    \end{align*}
    Thus,
    \begin{align*}
      \frac{I_{out}}{I_{in}} \approx \frac{I_{M_1}}{I_{M_0}} \approx \frac{k_{n_1}}{k_{n_0}} = \frac{\frac{\mu_n C_{ox} W_1}{L_1}}{\frac{\mu_n C_{ox} W_0}{L_0}} = \frac{\frac{W_1}{L_1}}{\frac{W_0}{L_0}} = \frac{W_1 L_0}{L_1 W_0}
    \end{align*}
  \item \textbf{Small Signal Analysis}:\\
    Applying KCL at $V_{out}$:
    \begin{align*}
      V_{out}G_3 + V_{out}g_{o_1} + g_{m_1}V_{x} &= 0 \\
      V_{out}(G_3 + g_{o_1}) + g_{m_1}V_{x} &= 0 \tag{5}
    \end{align*}
    Applying KCL at $V_{x}$:
    \begin{align*}
      V_{x}G_2 + V_{x}g_{o_0} + g_{m_0}V_{x} + I_{G} &= 0 \\
      V_{x}(G_2 + g_{o_0} + g_{m_0}) &= 0 \quad \text{[assuming } I_G \approx 0 \text{ by MOS physics]} \tag{6}
    \end{align*}
    Adding voltage source $V_s$ at $V_x$ that supplies
    a current $I_s$ into $V_{x}$, we get:
    \begin{align*}
      V_{x}(G_2 + g_{o_0} + g_{m_0}) - I_{s} &=0 \tag{7} \\
      V_{x} = V_{s} \tag{8}
    \end{align*}
    Solving equations (5), (7) and (8):
    \begin{align*}
      V_{out} &= - \frac{g_{m_1} V_{x}}{G_3 + g_{o_1}} \\
              &= - \frac{g_{m_1} V_{s}}{G_3 + g_{o_1}} \quad \text{[from (8)]} \\
      H(s) = \frac{N(s)}{D(s)} = \frac{V_{out}}{V_{s}} = - \frac{g_{m_1}}{G_3 + g_{o_1}}
    \end{align*}
\end{enumerate}

\subsubsection{Analysis}
\begin{enumerate}
  \item \textbf{Output Load (Max)}: Maximum possible
    load that can be connected at the output before
    the transistor $M_2$ goes out of saturation can
    be calculated as follows:
    \begin{align*}
      R_{out, \max} &= \frac{V_{L, \max}}{I_{M_1}} \quad \text{[where, } V_{ee} \leq V_{L, \max} \leq V_{cc} \text{]}
    \end{align*}
  \item \textbf{Impedances}:
    \paragraph{At Input:} Applying a test current $I_{test}$
      at $V_{x}$, from equation (6):
      \begin{align*}
        V_{x}(G_2 + g_{o_0} + g_{m_0}) &= I_{test} \\
        \therefore Z_{in} &= \frac{V_{x}}{I_{test}} = \frac{1}{G_2 + g_{o_0} + g_{m_0}}
      \end{align*}
      Ignoring $G_2$ and assuming $g_{m_0} \gg g_{o_0}$, we get:
      \begin{align*}
        Z_{in} &\approx \frac{1}{g_{m_0}} = \frac{2 I_{M_0}}{V_{ov_0}}
      \end{align*}
    \paragraph{At Output:} Applying a test current $I_{test}$
      at $V_{out}$, from equation (5):
      \begin{align*}
        V_{out}(G_3 + g_{o_1}) + g_{m_1}V_{x} &= I_{test} \\
        V_{out}(G_3 + g_{o_1}) + g_{m_1}(0) &= I_{test} \quad \text{[from (7)]} \\
        \therefore Z_{out} &= \frac{V_{out}}{I_{test}} = \frac{1}{G_3 + g_{o_1}}
      \end{align*}
      Ignoring $G_3$, we get:
      \begin{align*}
        Z_{out} &\approx \frac{1}{g_{o_1}} = \frac{1}{\lambda_{M_1} I_{M_1}}
      \end{align*}
\end{enumerate}

\subsubsection{Design}
Large Signal Analysis can be used to set the bias
currents and voltages in this circuit.
\begin{enumerate}
  \item Choose appropriate transistor dimensions
    ($W$ and $L$) for transistors $M_0$ and $M_1$
    to achieve the desired current gain.
    \begin{align*}
      \frac{I_{out}}{I_{in}} &= \frac{W_1 L_0}{L_1 W_0}
    \end{align*}
  \item Choose a suitable overdrive voltage value
    ($V_{ov}$) for computing the bias voltage at
    node $V_{x}$.
    \begin{align*}
      V_{ov_0} &= V_{gs_0} - V_{th} \\
      V_{ov_0} &= (V_{x} - V_{ee}) - V_{th} \\
      V_{x} &= V_{ov_0} + V_{th} + V_{ee}
    \end{align*}
  \item Using the chosen $V_{ov}$, compute $W_0$
    to set the required bias current $I_{in}$.
    \begin{align*}
      I_{in} = I_{M_0} &= \frac{1}{2} \mu_n C_{ox} \frac{W_0}{L_0} V_{ov_0}^2 \\
      \therefore W_0 &= \frac{2 I_{in} L_0}{\mu_n C_{ox} V_{ov_0}^2}
    \end{align*}
    $W_1$ can then be computed using the current
    gain relation from step (a).
  \item Compute the required value of $G_2$ (or
    $R_2$) to ensure that the voltage at node
    $V_{x}$ is as calculated in step (b).
    \begin{align*}
      (V_{x} - V_{cc})G_2 + I_{M_0} + I_{G} &= 0 \\
      (V_{x} - V_{cc})G_2 + I_{M_0} &\approx 0 \quad \text{[assuming } I_G \approx 0 \text{ by MOS physics ]} \\
      (V_{x} - V_{cc})G_2 &\approx - I_{M_0} \\
      G_2 &\approx - \frac{I_{M_0}}{(V_{x} - V_{cc})} \\
      \therefore G_2 &= \frac{I_{M_0}}{V_{cc} - V_{x}} \\
      R_2 &= \frac{1}{G_2} = \frac{V_{cc} - V_{x}}{I_{M_0}}
    \end{align*}
\end{enumerate}`

\subsubsection{Question}
Design a Widlar Current Mirror (Sink) that meets
the following requirements. The mirror must provide
a current gain at the output of 5. The input current
must be $I_{in} = 100 \mu A$. The mirror must operate
with $V_{cc} = 5V$ and $V_{ee} = 0V$. Assume the
following transistor parameters: $V_{th} = 1V$,
$\lambda = 0.02 V^{-1}$, $\frac{1}{2} \mu_n C_{ox} =
100 \mu A/V^2$ and $L = 1 \mu m$ for all transistors.
Find the output impedance of the mirror and the
maximum load that can be connected at the output
if $V_{L, \max} = 4V$.
\paragraph{Solution:}
\begin{enumerate}
  \item \textbf{Choose appropriate transistor dimensions:}
    \begin{align*}
      \frac{I_{out}}{I_{in}} &= \frac{W_1 L_0}{L_1 W_0} = 5 \\
      \therefore W_1 &= 5 W_0 \quad \text{[assuming } L_0 = L_1 \text{]}
    \end{align*}
  \item \textbf{Choose a suitable overdrive voltage and determine
    the resulting bias voltages:}
    \begin{align*}
      V_{x} &= V_{ov_0} + V_{th} + V_{ee} = 0.2V + 1V + 0V = 1.2V
    \end{align*}
  \item \textbf{Compute $W_0$ to set the required bias current:} For $I_{in} = 100 \mu A$
    \begin{align*}
      W_0 &= \frac{2 I_{in} L_0}{\mu_n C_{ox} V_{ov_0}^2} \\
          &= \frac{2 \cdot 100 \mu A \cdot 1 \mu m}{100 \mu A/V^2 \cdot (0.2V)^2} \\
          &= 25 \mu m \\
      \therefore W_1 &= 5 W_0 = 125 \mu m
    \end{align*}
  \item \textbf{Compute the required value of $G_2$ (or $R_2$):}
    \begin{align*}
      R_2 &= \frac{1}{G_2} = \frac{V_{cc} - V_{x}}{I_{M_0}} \\
          &= \frac{5V - 1.2V}{100 \mu A} \\
          &= 38 k\Omega
    \end{align*}
  \item \textbf{Compute the output impedance:}
    \begin{align*}
      Z_{out} &\approx \frac{1}{g_{o_1}} \approx \frac{1}{\lambda_{M_1} I_{M_1}} \\
              &= \frac{1}{0.02 V^{-1} \cdot 500 \mu A} \\
              &= 100 k\Omega
    \end{align*}
  \item \textbf{Compute the maximum load that can be connected
    at the output:}
    \begin{align*}
      R_{out, \max} &= \frac{V_{L, \max}}{I_{M_1}} \\
                   &= \frac{4V}{500 \mu A} \\
                   &= 8 k\Omega
    \end{align*}
\end{enumerate}

\subsubsection{Code}
\begin{lstlisting}[language=Octave]
  clear; clc;
  pkg load symbolic;

  syms Vx Vout Vs Is Itest
  syms G3 G2 go0 gm0 go1 gm1

  % system kcl
  eq1 = Vout*(G3 + go1) + Vx*gm1 == 0; % KCL @ Vout
  eq2 = Vx*(G2 + go0 + gm0) == 0; % KCL @ Vx

  % voltage transfer function: Vs forced at Vx & Is injected into Vx
  eq2_vt = Vx*(G2 + go0 + gm0) - Is == 0;
  eq3_vt = Vx == Vs; % equation introduced.
  sol_vt = solve([eq1, eq2_vt, eq3_vt], [Vout, Vx, Vs]);
  volt_transfer = simplify(sol_vt.Vout/sol_vt.Vs);
  % volt_transfer = subs(volt_transfer, 'G2', 0);
  % volt_transfer = subs(volt_transfer, 'G3', 0);
  printf('=== Voltage Transfer Results ===\n');
  pretty(volt_transfer)
  printf('\n');

  % current transfer function: Is injected into Vx
  sol_ct = solve([eq1, eq2_vt], [Vout, Vx]);
  Iout = sol_ct.Vout * G3;
  curr_transfer = simplify(Iout/Is);
  printf('=== Current Transfer Results ===\n');
  pretty(curr_transfer)
  printf('\n');

  % input impedance: Itest injected into Vx
  eq2_zin = Vx*(G2 + go0 + gm0) + Itest == 0;
  sol_zin = solve([eq1, eq2_zin], [Vout, Vx]);
  zin = simplify(sol_zin.Vx/Itest);
  % zin = subs(zin, 'G2', 0);
  % zin = subs(zin, 'G3', 0);
  printf('=== Input Impedance Results ===\n');
  pretty(zin)
  printf('\n');

  % output impedance: Itest injected into Vout
  eq1_zout = Vout*(G3 + go1) + Vx*gm1 + Itest == 0;
  sol_zout = solve([eq1_zout, eq2], [Vout, Vx]);
  zout = simplify(sol_zout.Vout/Itest);
  % zout = subs(zout, 'G2', 0);
  % zout = subs(zout, 'G3', 0);
  printf('=== Output Impedance Results ===\n');
  pretty(zout)
  printf('\n');
\end{lstlisting}

\subsection{Wilson Mirror}

\subsubsection{Overview}
Constructed from three transistors, It improves
upon Widlar Mirror by providing a higher output
impedance which results in better current
source/sink performance.

\subsubsection{Schematic}
\begin{center}
  \begin{tikzpicture}
  	% Paths, nodes and wires:
  	\node[shape=rectangle, draw={rgb,255:red,255;green,0;blue,0}, line width=1pt, dash pattern={on 4pt off 2pt on 1pt off 2pt on 1pt off 2pt}, minimum width=1.465cm, minimum height=2.195cm] at (6.25, 7.135){};
  	\draw (3.23, 8.04) to[american voltage source, l_={$V_{cc}$}] (1.98, 8.04);
  	\node[nmos, arrowmos, xscale=-1](N1) at (3.5, 2.75){} node[anchor=east] at (N1.text){$M_0$};
  	\node[nmos, arrowmos](N2) at (6, 2.75){} node[anchor=west] at (N2.text){$M_1$};
  	\draw (3.5, 7.54) to[european resistor, l={$G_3$}] (3.5, 6.04);
  	\draw (6, 5.79) |- (6, 5.25);
  	\draw (4.75, 2.75) |- (6, 3.71);
  	\draw (5.02, 2.75) -- (4.48, 2.75);
  	\draw (3.5, 1.98) -| (3.5, 1.75) -| (6, 1.98);
  	\draw (4, 1) to[american voltage source, l={$V_{ee}$}] (2.75, 1);
  	\draw (4, 1) -| (4.75, 1.75);
  	\node[ground, rotate=-90] at (1.98, 8.04){};
  	\node[ground, rotate=-90] at (2.75, 1){};
  	\draw[dash pattern={on 0.4pt off 0.8pt}, -stealth] (6, 5.79) -- (6.5, 5.79);
  	\node[shape=rectangle, minimum width=0.965cm, minimum height=0.465cm](N3) at (7, 5.75){} node[anchor=center] at (N3.text){$V_{out}$};
  	\draw (6, 7.5) to[european resistor, l={\textcolor{rgb,255:red,255;green,0;blue,0}{$G_4$}}] (6, 6);
  	\draw (6, 5.79) -- (6, 6.02);
  	\draw (3.5, 7.54) |- (3.23, 8.04);
  	\draw (3.5, 8.04) -| (6, 7.5);
  	\draw (10.25, 10.75) -- (10.25, -0);
  	\node[shape=rectangle, minimum width=4.59cm, minimum height=1.09cm] at (5.187, 9.813){} node[anchor=north, align=center, text width=4.202cm, inner sep=6pt] at (5.187, 10.375){Large Signal Schematic\\(Sink)};
  	\node[shape=rectangle, minimum width=3.09cm, minimum height=2.465cm] at (8.687, 7){} node[anchor=north west, align=left, text width=2.702cm, inner sep=6pt] at (7.125, 8.25){\textcolor{rgb,255:red,255;green,0;blue,0}{\footnotesize Exists Only For Analysis. In  Implementation take output directly from V\_\{out\}}};
  	\node[shape=rectangle, minimum width=0.965cm, minimum height=0.465cm](N4) at (7, 3.71){} node[anchor=center] at (N4.text){$V_{x}$};
  	\draw[dash pattern={on 0.4pt off 0.8pt}, -stealth] (6, 3.71) |- (6.5, 3.71);
  	\draw (3.5, 4.5) -| (3.5, 6.04);
  	\draw (17.625, 2.562) to[american current source, l_={$g_{m_1}v_{gs}$}] (17.625, 1.062);
  	\draw (19.125, 2.562) to[european resistor, l_={$g_{o_1}$}] (19.125, 1.062);
  	\draw (17.625, 1.062) -- (19.125, 1.062);
  	\draw (19.125, 2.562) -- (17.625, 2.562);
  	\draw (17.625, 1.062) -- (16.375, 1.062);
  	\draw (16.375, 2.562) -- (15.625, 2.562);
  	\node[shape=rectangle, minimum width=0.465cm, minimum height=0.465cm] at (16.375, 2.312){} node[anchor=center, align=center, text width=0.077cm, inner sep=6pt] at (16.375, 2.312){\footnotesize G};
  	\node[shape=rectangle, minimum width=0.465cm, minimum height=0.465cm] at (16.375, 1.312){} node[anchor=center, align=center, text width=0.077cm, inner sep=6pt] at (16.375, 1.312){\footnotesize S};
  	\node[shape=rectangle, minimum width=0.465cm, minimum height=0.465cm] at (19.625, 2.312){} node[anchor=center, align=center, text width=0.077cm, inner sep=6pt] at (19.625, 2.312){\footnotesize D};
  	\draw (15.625, 2.562) -- (14.875, 2.562);
  	\draw (13.375, 2.562) to[american current source, l={$g_{m_0}v_{gs}$}] (13.375, 1.062);
  	\draw (11.875, 2.562) to[european resistor, l={$g_{o_0}$}] (11.875, 1.062);
  	\draw (13.375, 2.563) -- (11.875, 2.563);
  	\draw (13.375, 1.063) -- (11.875, 1.063);
  	\draw (14.875, 1.062) -- (13.375, 1.062);
  	\node[shape=rectangle, minimum width=0.465cm, minimum height=0.465cm] at (14.875, 2.312){} node[anchor=center, align=center, text width=0.077cm, inner sep=6pt] at (14.875, 2.312){\footnotesize G};
  	\node[shape=rectangle, minimum width=0.465cm, minimum height=0.465cm] at (14.875, 1.312){} node[anchor=center, align=center, text width=0.077cm, inner sep=6pt] at (14.875, 1.312){\footnotesize S};
  	\node[shape=rectangle, minimum width=0.465cm, minimum height=0.465cm] at (11.375, 2.312){} node[anchor=center, align=center, text width=0.077cm, inner sep=6pt] at (11.375, 2.312){\footnotesize D};
  	\draw (11.875, 7.563) to[european resistor, l={$G_3$}] (11.875, 6.063);
  	\draw (11.875, 6.063) -| (11.875, 5.562);
  	\draw (11.875, 7.563) -| (11.875, 8.062) -- (12.625, 8.062) -| (12.625, 7.812);
  	\draw (19.125, 7.563) to[european resistor, l={\textcolor{rgb,255:red,255;green,0;blue,0}{$G_4$}}] (19.125, 6.063);
  	\draw (19.125, 6.063) -| (19.125, 5.562);
  	\draw (19.125, 7.563) -| (19.125, 8.062) |- (18.375, 8.062) -| (18.375, 7.812);
  	\node[ground] at (14.875, 1.062){};
  	\node[ground] at (16.375, 1.062){};
  	\node[ground] at (12.625, 7.812){};
  	\node[ground] at (18.375, 7.812){};
  	\node[shape=rectangle, minimum width=4.215cm, minimum height=1.09cm] at (15.5, 9.813){} node[anchor=north, align=center, text width=3.827cm, inner sep=6pt] at (15.5, 10.375){Small Signal Schematic\\(Sink)};
  	\draw (6, 3.52) -| (6, 3.71);
  	\node[nmos, arrowmos](N5) at (6, 4.48){} node[anchor=west] at (N5.text){$M_2$};
  	\draw (5.02, 4.48) -| (3.5, 3.52);
  	\node[shape=rectangle, minimum width=0.965cm, minimum height=0.465cm](N6) at (2.5, 4.5){} node[anchor=center] at (N6.text){$V_{y}$};
  	\draw[dash pattern={on 0.4pt off 0.8pt}, -stealth] (3.5, 4.48) -- (3, 4.48);
  	\draw (17.625, 5.562) to[american current source, l_={$g_{m_2}v_{gs}$}] (17.625, 4.062);
  	\draw (19.125, 5.562) to[european resistor, l_={$g_{o_2}$}] (19.125, 4.062);
  	\draw (17.625, 4.062) -- (19.125, 4.062);
  	\draw (19.125, 5.562) -- (17.625, 5.562);
  	\draw (17.625, 4.062) -- (16.375, 4.062);
  	\draw (16.375, 5.562) -- (15.625, 5.562);
  	\node[shape=rectangle, minimum width=0.465cm, minimum height=0.465cm] at (16.375, 5.312){} node[anchor=center, align=center, text width=0.077cm, inner sep=6pt] at (16.375, 5.312){\footnotesize G};
  	\node[shape=rectangle, minimum width=0.465cm, minimum height=0.465cm] at (16.375, 4.312){} node[anchor=center, align=center, text width=0.077cm, inner sep=6pt] at (16.375, 4.312){\footnotesize S};
  	\draw (19.125, 4.062) -- (19.125, 2.562);
  	\draw (15.625, 2.562) -| (15.625, 3.313) -- (19.125, 3.313);
  	\node[circ] at (3.5, 8.04){};
  	\node[circ] at (3.5, 4.48){};
  	\node[circ] at (4.75, 2.75){};
  	\node[circ] at (6, 3.71){};
  	\node[circ] at (4.75, 1.75){};
  	\node[circ] at (15.625, 2.562){};
  	\node[circ] at (19.125, 3.313){};
  	\draw (11.875, 2.562) |- (15.625, 5.562);
  	\node[shape=rectangle, minimum width=0.965cm, minimum height=0.465cm](N7) at (10.875, 5.582){} node[anchor=center] at (N7.text){$V_{y}$};
  	\draw[dash pattern={on 0.4pt off 0.8pt}, -stealth] (11.875, 5.562) |- (11.375, 5.562);
  	\node[circ] at (11.875, 5.562){};
  	\draw[dash pattern={on 0.4pt off 0.8pt}, -stealth] (19.125, 5.562) -- (19.625, 5.562);
  	\node[shape=rectangle, minimum width=0.965cm, minimum height=0.465cm](N8) at (20.125, 5.522){} node[anchor=center] at (N8.text){$V_{out}$};
  	\node[shape=rectangle, minimum width=0.965cm, minimum height=0.465cm](N9) at (14.625, 3.333){} node[anchor=center] at (N9.text){$I_G$};
  	\node[shape=rectangle, minimum width=0.965cm, minimum height=0.465cm](N10) at (20.125, 3.313){} node[anchor=center] at (N10.text){$V_{x}$};
  	\draw[dash pattern={on 0.4pt off 0.8pt}, -stealth] (19.125, 3.313) |- (19.625, 3.313);
  	\node[shape=rectangle, minimum width=0.965cm, minimum height=0.465cm](N11) at (3.75, 3.73){} node[anchor=center] at (N11.text){$I_G$};
  	\node[shape=rectangle, draw={rgb,255:red,255;green,0;blue,0}, line width=1pt, dash pattern={on 4pt off 2pt on 1pt off 2pt on 1pt off 2pt}, minimum width=1.59cm, minimum height=2.527cm] at (19.312, 7.094){};
  	\draw[dash pattern={on 0.4pt off 0.8pt}, -stealth] (4.75, 3.71) -- (4.25, 3.71);
  	\draw[dash pattern={on 0.4pt off 0.8pt}, -stealth] (15.625, 3.313) -- (15.125, 3.313);
  \end{tikzpicture}
\end{center}

\subsubsection{Transfer Function}
Here's a table of the voltages at each terminal of the
transistors:\\
\begin{tabularx}{\linewidth}{|X|X|X|X|X|X|} \hline
  Transistor & Gate (G) & Source (S) & Drain (D) & $V_{gs}$ & $V_{ds}$ \\ \hline
  $M_0$ & $V_{x}$ & $V_{ee}$ & $V_{y}$ & $V_{x} - V_{ee}$ & $V_{y} - V_{ee}$ \\ \hline
  $M_1$ & $V_{x}$ & $V_{ee}$ & $V_{x}$ & $V_{x} - V_{ee}$ & $V_{x} - V_{ee}$ \\ \hline
  $M_2$ & $V_{y}$ & $V_{x}$ & $V_{out}$ & $V_{y} - V_{x}$ & $V_{out} - V_{x}$ \\ \hline
\end{tabularx}\\
\begin{enumerate}
  \item \textbf{Large Signal Analysis}:\\
    Applying KCL at $V_{out}$:
    \begin{align*}
      (V_{out} - V_{cc})G_4 + I_{M_2} &= 0 \\
      V_{out} &= V_{cc} - \frac{I_{M_2}}{G_4} \tag{1}
    \end{align*}
    Applying KCL at $V_{y}$:
    \begin{align*}
      (V_{y} - V_{cc})G_3 + I_{M_0} &= 0 \\
      V_{y} &= V_{cc} - \frac{I_{M_0}}{G_3} \tag{2}
    \end{align*}
    Applying KCL at $V_{x}$:
    \begin{align*}
      I_{M_1} + I_{G} - I_{M_2} &= 0 \\
      I_{M_1} + I_{G} &= I_{M_2} \tag{3}
    \end{align*}
    Using transistor equations for $I_{M_0}$, $I_{M_1}$ and $I_{M_2}$
    and applying the voltages:
    \begin{align*}
      I_{M_0} &= k_{n_0}(V_{gs_0} - V_{th})^2 \cdot (1 + \lambda V_{ds_0}) \\
              &= k_{n_0}((V_{x} - V_{ee}) - V_{th})^2 \cdot (1 + \lambda (V_{y} - V_{ee})) \\
              &= k_{n_0}\left(V_{x} - V_{ee} - V_{th}\right)^2 \cdot \left(1 + \lambda \left(V_{cc} - \frac{I_{M_0}}{G_3} - V_{ee}\right)\right) \\
              &= k_{n_0}\left(V_{x} - V_{ee} - V_{th}\right)^2 \quad \text{[assuming } \lambda = 0 \text{]} \tag{4} \\ \\
      I_{M_1} &= k_{n_1}(V_{gs_1} - V_{th})^2 \cdot (1 + \lambda V_{ds_1}) \\
              &= k_{n_1}((V_{x} - V_{ee}) - V_{th})^2 \cdot (1 + \lambda (V_{x} - V_{ee})) \\
              &= k_{n_1}\left(V_{x} - V_{ee} - V_{th}\right)^2 \cdot \left(1 + \lambda (V_{x} - V_{ee})\right) \\
              &= k_{n_1}\left(V_{x} - V_{ee} - V_{th}\right)^2 \quad \text{[assuming } \lambda = 0 \text{]} \tag{5} \\ \\
      I_{M_2} &= k_{n_2}(V_{gs_2} - V_{th})^2 \cdot (1 + \lambda V_{ds_2}) \\
              &= k_{n_2}((V_{y} - V_{x}) - V_{th})^2 \cdot (1 + \lambda (V_{out} - V_{x})) \\
              &= k_{n_2}\left(V_{y} - V_{x} - V_{th}\right)^2 \cdot \left(1 + \lambda \left(V_{cc} - \frac{I_{M_2}}{G_4} - V_{x}\right)\right) \\
              &= k_{n_2}\left(V_{y} - V_{x} - V_{th}\right)^2 \quad \text{[assuming } \lambda = 0 \text{]} \tag{6} \\
    \end{align*}
    It can be seen that,
    \begin{align*}
      I_{in}  = I_{M_0} &\approx I_{M_1}  \quad \text{[if } k_{n_0} = k_{n_1} \text{]} \\
      I_{out} = I_{M_2} &= I_{M_1} + I_{G} \approx I_{M_1} \quad \text{[as } I_{G} \approx 0 \text{ by MOS physics]}
    \end{align*}
    Thus,
    \begin{align*}
      \boxed{\frac{I_{out}}{I_{in}} \approx \frac{I_{M_2}}{I_{M_0}} \approx \frac{I_{M_1}}{I_{M_0}} \approx \frac{k_{n_1}}{k_{n_0}} = \frac{\frac{\mu_n C_{ox} W_1}{L_1}}{\frac{\mu_n C_{ox} W_0}{L_0}} = \frac{\frac{W_1}{L_1}}{\frac{W_0}{L_0}} = \frac{W_1 L_0}{L_1 W_0}}
    \end{align*}
  \item \textbf{Small Signal Analysis}:\\
    Applying KCL at $V_{out}$:
    \begin{align*}
      V_{out}G_4 + (V_{out} - V_{x})g_{o_2} + g_{m_2}(V_{y} - V_{x}) &= 0 \\
      V_{out}(G_4 + g_{o_2}) - V_{x}(g_{o_2} + g_{m_2}) + g_{m_2}V_{y} &= 0 \tag{7} \\
    \end{align*}
    Applying KCL at $V_{x}$:
    \begin{align*}
      (V_{x} - V_{out})g_{o_2} + V_{x}g_{o_1} + g_{m_1}V_{x} - g_{m_2}(V_{y} - V_{x}) + I_{G} &= 0 \\
      (V_{x} - V_{out})g_{o_2} + V_{x}g_{o_1} + g_{m_1}V_{x} - g_{m_2}(V_{y} - V_{x}) &= 0 \quad \text{[as } I_{G} \approx 0 \text{ by MOS physics]} \\
      V_{x}g_{o_2} - V_{out}g_{o_2} + V_{x}g_{o_1} + g_{m_1}V_{x} - g_{m_2}V_{y} + g_{m_2}V_{x} &= 0 \\
      V_{x}(g_{o_2} + g_{o_1} + g_{m_1} + g_{m_2}) - V_{out}g_{o_2} - g_{m_2}V_{y} &= 0 \tag{8}
    \end{align*}
    Applying KCL at $V_{y}$:
    \begin{align*}
      V_{y}G_3 + V_{y}g_{o_0} + g_{m_0}V_{x} &= 0 \\
      V_{y}(G_3 + g_{o_0}) + g_{m_0}V_{x} &= 0 \tag{9}
    \end{align*}
    Adding voltage source $V_s$ at $V_y$ that supplies
    a current $I_s$ into $V_y$, we get:
    \begin{align*}
      V_{y}(G_3 + g_{o_0}) + g_{m_0}V_{x} - I_{s} &= 0 \tag{10}
    \end{align*}
    Fourth equation arises from Voltage source:
    \begin{align*}
      V_{y} &= V_{s} \tag{11}
    \end{align*}
    Solving equations (7), (8), (10) and (11),
    \begin{align*}
      V_{out}(G_4 + g_{o_2}) - V_{x}(g_{o_2} + g_{m_2}) + g_{m_2}V_{y} &= 0 \\
      V_{out}(G_4 + g_{o_2}) - V_{x}(g_{o_2} + g_{m_2}) + g_{m_2}V_{s} &= 0 \quad \text{[from (11)]} \tag{12} \\
      \\
      V_{x}(g_{o_2} + g_{o_1} + g_{m_1} + g_{m_2}) - V_{out}g_{o_2} - g_{m_2}V_{y} &= 0 \\
      V_{x}(g_{o_2} + g_{o_1} + g_{m_1} + g_{m_2}) - V_{out}g_{o_2} - g_{m_2}V_{s} &= 0 \quad \text{[from (11)]} \\
      V_{x}(g_{o_2} + g_{o_1} + g_{m_1} + g_{m_2}) &= V_{out}g_{o_2} + g_{m_2}V_{s} \\
      V_{x} &= \frac{V_{out}g_{o_2} + g_{m_2}V_{s}}{g_{o_2} + g_{o_1} + g_{m_1} + g_{m_2}} \tag{13} \\
    \end{align*}
    Putting (13) in (12),
    \begin{align*}
      V_{out}(G_4 + g_{o_2}) - \left(\frac{V_{out}g_{o_2} + g_{m_2}V_{s}}{g_{o_2} + g_{o_1} + g_{m_1} + g_{m_2}}\right)(g_{o_2} + g_{m_2}) + g_{m_2}V_{s} &= 0 \\
      V_{out}(G_4 + g_{o_2}) - \frac{V_{out}g_{o_2}(g_{o_2} + g_{m_2})}{g_{o_2} + g_{o_1} + g_{m_1} + g_{m_2}} - \frac{g_{m_2}V_{s}(g_{o_2} + g_{m_2})}{g_{o_2} + g_{o_1} + g_{m_1} + g_{m_2}} + g_{m_2}V_{s} &= 0 \\
      V_{out}\left((G_4 + g_{o_2}) - \frac{g_{o_2}(g_{o_2} + g_{m_2})}{g_{o_2} + g_{o_1} + g_{m_1} + g_{m_2}}\right) + V_{s}\left(g_{m_2} - \frac{g_{m_2}(g_{o_2} + g_{m_2})}{g_{o_2} + g_{o_1} + g_{m_1} + g_{m_2}}\right) &= 0 \\
      V_{out}\left(\frac{(G_4 + g_{o_2})(g_{o_2} + g_{o_1} + g_{m_1} + g_{m_2}) - g_{o_2}(g_{o_2} + g_{m_2})}{g_{o_2} + g_{o_1} + g_{m_1} + g_{m_2}}\right) \\ \qquad \, = -V_{s}\left(\frac{g_{m_2}(g_{o_2} + g_{o_1} + g_{m_1} + g_{m_2}) - g_{m_2}(g_{o_2} + g_{m_2})}{g_{o_2} + g_{o_1} + g_{m_1} + g_{m_2}}\right) \\
      V_{out}\left(\frac{G_4(g_{o_2} + g_{o_1} + g_{m_1} + g_{m_2}) + g_{o_2}(g_{o_2} + g_{o_1} + g_{m_1} + g_{m_2}) - g_{o_2}(g_{o_2} + g_{m_2})}{g_{o_2} + g_{o_1} + g_{m_1} + g_{m_2}}\right) \\ \qquad \, = -V_{s}\left(\frac{g_{m_2}(g_{o_2} + g_{o_1} + g_{m_1} + g_{m_2}) - g_{m_2}(g_{o_2} + g_{m_2})}{g_{o_2} + g_{o_1} + g_{m_1} + g_{m_2}}\right) \\
      V_{out}\left(\frac{G_4(g_{o_2} + g_{o_1} + g_{m_1} + g_{m_2}) + g_{o_2}(g_{o_1} + g_{m_1}) + g_{o_2}(g_{o_2} + g_{m_2}) - g_{o_2}(g_{o_2} + g_{m_2})}{g_{o_2} + g_{o_1} + g_{m_1} + g_{m_2}}\right) \\ \qquad \, = -V_{s}\left(\frac{g_{m_2}(g_{o_2} + g_{o_1} + g_{m_1} + g_{m_2}) - g_{m_2}(g_{o_2} + g_{m_2})}{g_{o_2} + g_{o_1} + g_{m_1} + g_{m_2}}\right) \\
      V_{out}\left(\frac{G_4(g_{o_2} + g_{o_1} + g_{m_1} + g_{m_2}) + g_{o_2}(g_{o_1} + g_{m_1})}{g_{o_2} + g_{o_1} + g_{m_1} + g_{m_2}}\right) \\ \qquad \, = -V_{s}\left(\frac{g_{m_2}(g_{o_2} + g_{o_1} + g_{m_1} + g_{m_2}) - g_{m_2}(g_{o_2} + g_{m_2})}{g_{o_2} + g_{o_1} + g_{m_1} + g_{m_2}}\right) \\
      V_{out}\left(\frac{G_4(g_{o_2} + g_{o_1} + g_{m_1} + g_{m_2}) + g_{o_2}(g_{o_1} + g_{m_1})}{g_{o_2} + g_{o_1} + g_{m_1} + g_{m_2}}\right) \\ \qquad \, = -V_{s}\left(\frac{g_{m_2}(g_{o_2} + g_{o_1} + g_{m_1} + g_{m_2}) - g_{m_2}(g_{o_2} + g_{m_2})}{g_{o_2} + g_{o_1} + g_{m_1} + g_{m_2}}\right) \\
      V_{out}(G_4(g_{o_2} + g_{o_1} + g_{m_1} + g_{m_2}) + g_{o_2}(g_{o_1} + g_{m_1})) \\ \qquad \, = -V_{s}(g_{m_2}(g_{o_2} + g_{o_1} + g_{m_1} + g_{m_2}) - g_{m_2}(g_{o_2} + g_{m_2})) \\\\
      V_{out}(G_4(g_{o_2} + g_{o_1} + g_{m_1} + g_{m_2}) + g_{o_2}(g_{o_1} + g_{m_1})) \\ \qquad \, = -V_{s}(g_{m_2}g_{o_2} + g_{m_2}g_{o_1} + g_{m_2}g_{m_1} + g_{m_2}^2 - g_{m_2}g_{o_2} - g_{m_2}^2) \\\\
      V_{out}(G_4(g_{o_2} + g_{o_1} + g_{m_1} + g_{m_2}) + g_{o_2}(g_{o_1} + g_{m_1})) \\ \qquad \, = -V_{s}(g_{m_2}g_{o_1} + g_{m_2}g_{m_1}) \\\\
      V_{out}(G_4(g_{o_2} + g_{o_1} + g_{m_1} + g_{m_2}) + g_{o_2}(g_{o_1} + g_{m_1})) \\ \qquad \, = -V_{s}(g_{m_2}(g_{o_1} + g_{m_1})) \\\\
      \therefore \frac{V_{out}}{V_s} = \frac{-g_{m_2}(g_{o_1} + g_{m_1})}{G_4(g_{o_2} + g_{o_1} + g_{m_1} + g_{m_2}) + g_{o_2}(g_{o_1} + g_{m_1})} \tag{14} \\
    \end{align*}
\end{enumerate}

\subsubsection{Analysis}
\begin{enumerate}
  \item \textbf{Output Load (Max)}: Maximum possible
    load that can be connected at the output before
    the transistor $M_2$ goes out of saturation can
    be calculated as follows:
    \begin{align*}
      R_{out, max} &= \frac{V_{L, max}}{I_{M_2}} \quad \text{[where, } V_{ee} \leq V_{L, max} \leq V_{cc} \text{]}
    \end{align*}
  \item \textbf{Impedance}:
    \paragraph{At Input:} Applying a test current $I_{test}$
      at $V_{y}$, from equation (9):
      \begin{align*}
        V_{y}(G_3 + g_{o_0}) + g_{m_0}V_{x} &= I_{test} \\
      \end{align*}
      Substituting (14) into (13):
      \begin{align*}
        V_{x} &= \frac{V_{out}g_{o_2} + g_{m_2}V_{y}}{g_{o_2} + g_{o_1} + g_{m_1} + g_{m_2}} \\
              &= \frac{\left(V_{y} \cdot \frac{-g_{m_2}(g_{o_1} + g_{m_1})}{G_4(g_{o_2} + g_{o_1} + g_{m_1} + g_{m_2}) + g_{o_2}(g_{o_1} + g_{m_1})}\right)g_{o_2} + g_{m_2}V_{y}}{g_{o_2} + g_{o_1} + g_{m_1} + g_{m_2}} \\
              &= \frac{V_{y}\frac{-g_{m_2}g_{o_2}(g_{o_1} + g_{m_1})}{G_4(g_{o_2} + g_{o_1} + g_{m_1} + g_{m_2}) + g_{o_2}(g_{o_1} + g_{m_1})} + g_{m_2}V_{y}}{g_{o_2} + g_{o_1} + g_{m_1} + g_{m_2}} \\
              &= V_{y} \cdot \frac{g_{m_2} - \frac{g_{m_2}g_{o_2}(g_{o_1} + g_{m_1})}{G_4(g_{o_2} + g_{o_1} + g_{m_1} + g_{m_2}) + g_{o_2}(g_{o_1} + g_{m_1})}}{g_{o_2} + g_{o_1} + g_{m_1} + g_{m_2}} \\
              &= V_{y} \cdot \frac{g_{m_2}G_4(g_{o_2} + g_{o_1} + g_{m_1} + g_{m_2}) + g_{m_2}g_{o_2}(g_{o_1} + g_{m_1}) - g_{m_2}g_{o_2}(g_{o_1} + g_{m_1})}{(g_{o_2} + g_{o_1} + g_{m_1} + g_{m_2})G_4(g_{o_2} + g_{o_1} + g_{m_1} + g_{m_2}) + g_{o_2}(g_{o_1} + g_{m_1})(g_{o_2} + g_{o_1} + g_{m_1} + g_{m_2})} \\
              &= V_{y} \cdot \frac{g_{m_2}G_4(g_{o_2} + g_{o_1} + g_{m_1} + g_{m_2})}{(g_{o_2} + g_{o_1} + g_{m_1} + g_{m_2})G_4(g_{o_2} + g_{o_1} + g_{m_1} + g_{m_2}) + g_{o_2}(g_{o_1} + g_{m_1})(g_{o_2} + g_{o_1} + g_{m_1} + g_{m_2})} \\
              &= V_{y} \cdot \frac{g_{m_2}G_4}{G_4(g_{o_2} + g_{o_1} + g_{m_1} + g_{m_2}) + g_{o_2}(g_{o_1} + g_{m_1})} \tag{15} \\
      \end{align*}
      Substituting (15) into (9):
      \begin{align*}
        V_{y}(G_3 + g_{o_0}) + g_{m_0}\left(V_{y} \cdot \frac{g_{m_2}G_4}{G_4(g_{o_2} + g_{o_1} + g_{m_1} + g_{m_2}) + g_{o_2}(g_{o_1} + g_{m_1})}\right) &= I_{test} \\
        V_{y}\left((G_3 + g_{o_0}) + \frac{g_{m_0}g_{m_2}G_4}{G_4(g_{o_2} + g_{o_1} + g_{m_1} + g_{m_2}) + g_{o_2}(g_{o_1} + g_{m_1})}\right) &= I_{test} \\
        V_{y}\left(\frac{(G_3 + g_{o_0})(G_4(g_{o_2} + g_{o_1} + g_{m_1} + g_{m_2}) + g_{o_2}(g_{o_1} + g_{m_1})) + g_{m_0}g_{m_2}G_4}{G_4(g_{o_2} + g_{o_1} + g_{m_1} + g_{m_2}) + g_{o_2}(g_{o_1} + g_{m_1})}\right) &= I_{test} \\
        \therefore Z_{in} = \frac{V_{y}}{I_{test}} = \frac{G_4(g_{o_2} + g_{o_1} + g_{m_1} + g_{m_2}) + g_{o_2}(g_{o_1} + g_{m_1})}{(G_3 + g_{o_0})(G_4(g_{o_2} + g_{o_1} + g_{m_1} + g_{m_2}) + g_{o_2}(g_{o_1} + g_{m_1})) + g_{m_0}g_{m_2}G_4} \\
      \end{align*}
    \paragraph{At Output:} Applying a test current $I_{test}$
      at $V_{out}$, from equation (7):
      \begin{align*}
        V_{out}(G_4 + g_{o_2}) - V_{x}(g_{o_2} + g_{m_2}) + g_{m_2}V_{y} &= I_{test} \\
        V_{out}(G_4 + g_{o_2}) - V_{x}(g_{o_2} + g_{m_2}) + g_{m_2}\left(\frac{-g_{m_0}V_{x}}{G_3 + g_{o_0}}\right) &= I_{test} \quad \text{[rearraning (9)]} \\
        V_{out}(G_4 + g_{o_2}) - V_{x}\left((g_{o_2} + g_{m_2}) + \frac{g_{m_2}g_{m_0}}{G_3 + g_{o_0}}\right) &= I_{test} \\
        V_{out}(G_4 + g_{o_2}) - V_{x}\left(\frac{(g_{o_2} + g_{m_2})(G_3 + g_{o_0}) + g_{m_2}g_{m_0}}{G_3 + g_{o_0}}\right) &= I_{test} \tag{16}
      \end{align*}
      Substituting (14) into (13):
      \begin{align*}
        V_{x} &= \frac{V_{out}g_{o_2} + g_{m_2}V_{y}}{g_{o_2} + g_{o_1} + g_{m_1} + g_{m_2}} \\
              &= \frac{V_{out}g_{o_2} + g_{m_2}\left(V_{out}\left(\frac{G_4(g_{o_2} + g_{o_1} + g_{m_1} + g_{m_2}) + g_{o_2}(g_{o_1} + g_{m_1})}{-g_{m_2}(g_{o_1} + g_{m_1})}\right)\right)}{g_{o_2} + g_{o_1} + g_{m_1} + g_{m_2}} \\
              &= \frac{V_{out}g_{o_2} - V_{out}\left(\frac{G_4(g_{o_2} + g_{o_1} + g_{m_1} + g_{m_2}) + g_{o_2}(g_{o_1} + g_{m_1})}{(g_{o_1} + g_{m_1})}\right)}{g_{o_2} + g_{o_1} + g_{m_1} + g_{m_2}} \\
              &= V_{out} \left(\frac{g_{o_2} - \left(\frac{G_4(g_{o_2} + g_{o_1} + g_{m_1} + g_{m_2}) + g_{o_2}(g_{o_1} + g_{m_1})}{(g_{o_1} + g_{m_1})}\right)}{g_{o_2} + g_{o_1} + g_{m_1} + g_{m_2}}\right) \\
              &= V_{out} \left(\frac{\frac{g_{o_2}(g_{o_1} + g_{m_1}) - G_4(g_{o_2} + g_{o_1} + g_{m_1} + g_{m_2}) - g_{o_2}(g_{o_1} + g_{m_1})}{(g_{o_1} + g_{m_1})}}{g_{o_2} + g_{o_1} + g_{m_1} + g_{m_2}}\right) \\
              &= V_{out} \left(\frac{- G_4(g_{o_2} + g_{o_1} + g_{m_1} + g_{m_2})}{(g_{o_1} + g_{m_1})(g_{o_2} + g_{o_1} + g_{m_1} + g_{m_2})}\right) \\
              &= V_{out} \left(\frac{- G_4}{(g_{o_1} + g_{m_1})}\right) \tag{17}
      \end{align*}`
      Substituting (17) into (16):
      \begin{align*}
        V_{out}(G_4 + g_{o_2}) - \left(V_{out} \left(\frac{-G_4}{(g_{o_1} + g_{m_1})}\right)\right)\left(\frac{(g_{o_2} + g_{m_2})(G_3 + g_{o_0}) + g_{m_2}g_{m_0}}{G_3 + g_{o_0}}\right) &= I_{test} \\
        V_{out}(G_4 + g_{o_2}) + \left(V_{out} \cdot \frac{G_4(g_{o_2} + g_{m_2})(G_3 + g_{o_0}) + G_4g_{m_2}g_{m_0}}{(g_{o_1} + g_{m_1})(G_3 + g_{o_0})} \right) &= I_{test} \\
        V_{out}\left((G_4 + g_{o_2}) + \left(\frac{G_4(g_{o_2} + g_{m_2})(G_3 + g_{o_0}) + G_4g_{m_2}g_{m_0}}{(g_{o_1} + g_{m_1})(G_3 + g_{o_0})}\right)\right) &= I_{test} \\
        V_{out}\left(\frac{(G_4 + g_{o_2})(g_{o_1} + g_{m_1})(G_3 + g_{o_0}) + G_4(g_{o_2} + g_{m_2})(G_3 + g_{o_0}) + G_4g_{m_2}g_{m_0}}{(g_{o_1} + g_{m_1})(G_3 + g_{o_0})}\right) &= I_{test} \\
        V_{out}\left(\frac{(G_3 + g_{o_0})((G_4 + g_{o_2})(g_{o_1} + g_{m_1}) + G_4(g_{o_2} + g_{m_2})) + G_4g_{m_2}g_{m_0}}{(g_{o_1} + g_{m_1})(G_3 + g_{o_0})}\right) &= I_{test} \\
        V_{out}\left(\frac{(G_3 + g_{o_0})(G_4(g_{o_1} + g_{m_1}) + g_{o_2}(g_{o_1} + g_{m_1}) + G_4(g_{o_2} + g_{m_2})) + G_4g_{m_2}g_{m_0}}{(g_{o_1} + g_{m_1})(G_3 + g_{o_0})}\right) &= I_{test} \\
        V_{out}\left(\frac{(G_3 + g_{o_0})(G_4(g_{o_1} + g_{m_1} + g_{o_2} + g_{m_2}) + g_{o_2}(g_{o_1} + g_{m_1})) + G_4g_{m_2}g_{m_0}}{(g_{o_1} + g_{m_1})(G_3 + g_{o_0})}\right) &= I_{test} \\
        \therefore Z_{out} = \frac{V_{out}}{I_{test}} = \frac{(g_{o_1} + g_{m_1})(G_3 + g_{o_0})}{(G_3 + g_{o_0})(G_4(g_{o_1} + g_{m_1} + g_{o_2} + g_{m_2}) + g_{o_2}(g_{o_1} + g_{m_1})) + G_4g_{m_2}g_{m_0}}
      \end{align*}
\end{enumerate}

\subsubsection{Design}
Large Signal Analysis can be used to set the bias
currents and voltages in this circuit.
\begin{enumerate}
  \item Choose appropriate transistor dimensions
    ($W$ and $L$) for transistors $M_0$ and $M_1$
    to achieve the desired current gain.
    \begin{align*}
      \frac{I_{out}}{I_{in}} &= \frac{W_1 L_0}{L_1 W_0}
    \end{align*}
  \item Choose a suitable overdrive voltage value
    ($V_{ov}$) for setting the bias voltages at
    nodes $V_{x}$ and $V_{y}$.
    \begin{align*}
      V_{ov_0} &= V_{gs_0} - V_{th} \\
      V_{ov_0} &= (V_{x} - V_{ee}) - V_{th} \\
      V_{x} &= V_{ov_0} + V_{th} + V_{ee} \\
      \\
      V_{ov_2} &= V_{gs_2} - V_{th} \\
      V_{ov_2} &= (V_{y} - V_{x}) - V_{th} \\
      V_{y} &= V_{ov_2} + V_{th} + V_{x}
    \end{align*}
  \item Using the chosen $V_{ov}$, compute $W_0$
    to set the required bias current $I_{in}$.
    \begin{align*}
      I_{in} = I_{M_0} &= \frac{1}{2} \mu_n C_{ox} \frac{W_0}{L_0} V_{ov_0}^2 \\
      \therefore W_0 &= \frac{2 I_{in} L_0}{\mu_n C_{ox} V_{ov_0}^2}
    \end{align*}
    $W_1$ can then be computed using the current
    gain relation from step (a).
  \item Compute the required value of $G_3$ (or
    $R_3$) to ensure that the voltage at node
    $V_{y}$ is as calculated in step (b).
    \begin{align*}
      (V_{y} - V_{cc})G_3 + I_{M_0} &= 0 \\
      \therefore G_3 &= \frac{I_{in}}{V_{cc} - V_{y}} \quad \text{[}I_{M_0} = I_{in} \text{]} \\
      R_3 &= \frac{1}{G_3} = \frac{V_{cc} - V_{y}}{I_{in}}
    \end{align*}
\end{enumerate}

\subsubsection{Question}
Design a Wilson Current Mirror (Sink) that meets the
following requirements. The mirror must provide a current
gain at the output of 5. The input current must be
$I_{in} = 100 \mu A$. The mirror must operate with
$V_{cc} = 5V$ and $V_{ee} = 0V$. Assume the following
transistor parameters: $V_{th} = 1V$, $\lambda = 0.02 V^{-1}$,
$\frac{1}{2} \mu_n C_{ox} = 100 \mu A/V^2$ and
$L = 1 \mu m$ for all transistors. Find the output impedance
and the maximum possible load that can be connected at the
output assuming $V_{L, max} = V_{cc}$.
\paragraph{Solution:}
\begin{enumerate}
  \item \textbf{Choose appropriate transistor dimensions:}
    \begin{align*}
      \frac{I_{out}}{I_{in}} &= \frac{W_1 L_0}{L_1 W_0} = 5 \\
      \text{Assuming } L_0 &= L_1 = 1 \mu m \\
      \therefore W_1 &= 5 W_0 \tag{i}
    \end{align*}
  \item \textbf{Choose a suitable overdrive voltage and determine
    the resulting bias voltages:} Assuming $V_{ov} = 0.2V$,
    \begin{align*}
      V_{x} &= V_{ov_0} + V_{th} + V_{ee} = 0.2V + 1V + 0V = 1.2V \\
      V_{y} &= V_{ov_2} + V_{th} + V_{x} = 0.2V + 1V + 1.2V = 2.4V
    \end{align*}
  \item \textbf{Compute $W_0$ to set the required bias current:}
    For $I_{in} = 100 \mu A$
    \begin{align*}
      W_0 &= \frac{2 I_{in} L_0}{\mu_n C_{ox} V_{ov_0}^2} \\
          &= \frac{2 \cdot 100 \mu A \cdot 1 \mu m}{100 \mu A/V^2 \cdot (0.2V)^2} \\
          &= 25 \mu m \\
      W_1 &= 5 W_0 = 125 \mu m
    \end{align*}
  \item \textbf{Compute the required value of $G_3$ (or $R_3$):}
    \begin{align*}
      R_3 &= \frac{1}{G_3} = \frac{V_{cc} - V_{y}}{I_{in}} \\
          &= \frac{(5 - 2.4)}{100 \mu A} = 26k \Omega
    \end{align*}
  \item \textbf{Compute the output impedance:}
    \begin{align*}
      Z_{out} &\approx \frac{2}{\lambda_{M_1} \lambda_{M_2} I_{M_1} ((V_{y} - V_{x}) - V_{th})} \\
              &= \frac{2}{0.02 \cdot 0.02 \cdot 500 \mu A \cdot (2.4V - 1.2V - 1V)} \\
              &= 50 M \Omega
    \end{align*}
  \item \textbf{Compute the maximum load that can be connected
    at the output:}
    \begin{align*}
      R_{out, max} &= \frac{V_{L, max}}{I_{M_2}} \\
                   &= \frac{5V}{500 \mu A} \\
                   &= 10 k\Omega
    \end{align*}
\end{enumerate}

\subsubsection{Code}
\begin{lstlisting}[language=Octave]
  clear; clc;
  pkg load symbolic;

  syms Vy Vx Vout Vs Is Itest
  syms G3 G4 go0 gm0 go1 gm1 go2 gm2

  % system kcl
  eq1 = Vout*(G4 + go2) - Vx*(go2 + gm2) + gm2*Vy == 0; % KCL @ Vout
  eq2 = -Vout*go2 + Vx*(go2 + go1 + gm1 + gm2) - gm2*Vy == 0; % KCL @ Vx
  eq3 = Vy*(G3 + go0) + gm0*Vx == 0; % KCL @ Vy

  % voltage transfer function: Vs forced at Vy & Is injected into Vy
  eq3_vt = Vy*(G3 + go0) + gm0*Vx - Is == 0;
  eq4_vt = Vy == Vs;
  sol_vt = solve([eq1, eq2, eq3_vt, eq4_vt], [Vout, Vx, Vy, Vs]);
  volt_transfer = simplify(sol_vt.Vout/sol_vt.Vs);
  % volt_transfer = collect(volt_transfer, G4);
  % volt_transfer = subs(volt_transfer, 'G3', 0);
  % volt_transfer = subs(volt_transfer, 'G4', 0);
  printf('=== Voltage Transfer Results ===\n');
  pretty(volt_transfer)
  printf('\n');

  % current transfer function: Is injected into Vy
  sol_ct = solve([eq1, eq2, eq3_vt], [Vout, Vx, Vy]);
  Iout = sol_ct.Vout * G4;
  curr_transfer = simplify(Iout/Is);
  % curr_transfer = collect(curr_transfer, G3*G4);
  printf('=== Current Transfer Results ===\n');
  pretty(curr_transfer)
  printf('\n');

  % input impedence: Itest injected into Vy
  eq3_zin = Vy*(G3 + go0) + gm0*Vx - Itest == 0;
  sol_zin = solve([eq1, eq2, eq3_zin], [Vout, Vx, Vy]);
  zin = simplify(sol_zin.Vy/Itest);
  % zin = collect(zin, G3*G4);
  % zin = subs(zin, 'G3', 0);
  % zin = subs(zin, 'G4', 0);
  printf('=== Input Impedence Results ===\n');
  pretty(zin)
  printf('\n');

  % output impedence: Itest injected into Vout
  eq1_zout = Vout*(G4 + go2) - Vx*(go2 + gm2) + gm2*Vy - Itest == 0;
  sol_zout = solve([eq1_zout, eq2, eq3], [Vout, Vx, Vy]);
  zout = simplify(sol_zout.Vout/Itest);
  % zout = collect(zout, G3*G4);
  % zout = subs(zout, 'G3', 0);
  % zout = subs(zout, 'G4', 0);
  printf('=== Output Impedence Results ===\n');
  pretty(zout)
  printf('\n');
\end{lstlisting}

\subsection{Modified Wilson Mirror}

\paragraph{Overview:}
Constructed from transistors and resistors, a
Modified Wilson Mirror is an enhanced version of
the Wilson Current Mirror. It offers improved
performance in terms of output resistance and
accuracy of current replication.

\paragraph{Schematic:}
\begin{center}
  \begin{tikzpicture}
  	% Paths, nodes and wires:
  	\node[shape=rectangle, draw={rgb,255:red,255;green,0;blue,0}, line width=1pt, dash pattern={on 4pt off 2pt on 1pt off 2pt on 1pt off 2pt}, minimum width=1.965cm, minimum height=2.215cm] at (15.02, 2.688){};
  	\node[shape=rectangle, draw={rgb,255:red,255;green,0;blue,0}, line width=1pt, dash pattern={on 4pt off 2pt on 1pt off 2pt on 1pt off 2pt}, minimum width=1.965cm, minimum height=2.215cm] at (6, 7.145){};
  	\draw (3.23, 8.04) to[american voltage source, l_={$V_{cc}$}] (1.98, 8.04);
  	\node[nmos, arrowmos, xscale=-1](N1) at (3.5, 2.75){} node[anchor=east] at (N1.text){$M_0$};
  	\node[nmos, arrowmos](N2) at (6, 2.75){} node[anchor=west] at (N2.text){$M_1$};
  	\draw (3.5, 7.54) to[european resistor, l={$G_4$}] (3.5, 6.04);
  	\draw (6, 5.79) |- (6, 5.25);
  	\draw (4.75, 2.75) |- (6, 3.71);
  	\draw (5.02, 2.75) -- (4.48, 2.75);
  	\draw (3.5, 1.98) -| (3.5, 1.75) -| (6, 1.98);
  	\draw (4, 1) to[american voltage source, l={$V_{ee}$}] (2.75, 1);
  	\draw (4, 1) -| (4.75, 1.75);
  	\node[ground, rotate=-90] at (1.98, 8.04){};
  	\node[ground, rotate=-90] at (2.75, 1){};
  	\draw (6, 5.79) -- (6.5, 5.79);
  	\node[shape=rectangle, minimum width=0.965cm, minimum height=0.465cm](N3) at (7, 5.75){} node[anchor=center] at (N3.text){$V_{out}$};
  	\draw (6, 7.5) to[european resistor, l={\textcolor{rgb,255:red,255;green,0;blue,0}{$G_5$}}] (6, 6);
  	\draw (6, 5.79) -- (6, 6.02);
  	\draw (3.5, 7.54) |- (3.23, 8.04);
  	\path[draw={rgb,255:red,255;green,0;blue,0}] (3.5, 8.04) -| (6, 7.5);
  	\node[pmos, arrowmos](N4) at (15, 7.168){} node[anchor=west] at (N4.text){$M_1$};
  	\node[pmos, arrowmos, xscale=-1](N5) at (12.5, 7.168){} node[anchor=east] at (N5.text){$M_0$};
  	\draw (13.98, 7.168) -- (13.48, 7.168);
  	\draw (15, 4.71) -| (15, 4.27);
  	\draw (13.71, 7.168) |- (14.77, 6.25);
  	\draw (12.5, 3.75) to[european resistor, l={$G_4$}] (12.5, 2.25);
  	\draw (15, 3.75) to[european resistor, l={\textcolor{rgb,255:red,255;green,0;blue,0}{$G_5$}}] (15, 2.25);
  	\draw (15, 4.27) |- (15, 3.75);
  	\draw (12.5, 7.938) |- (15, 8.188) -| (15, 7.938);
  	\draw (15, 4.02) -- (15.5, 4.02);
  	\node[shape=rectangle, minimum width=0.965cm, minimum height=0.465cm](N6) at (16, 4.02){} node[anchor=center] at (N6.text){$V_{out}$};
  	\draw (12.5, 2.25) -- (12.5, 2) -| (13.75, 2);
  	\draw (12.5, 8.188) to[american voltage source, l_={$V_{cc}$}] (11.25, 8.188);
  	\node[ground, rotate=-90] at (11.25, 8.188){};
  	\draw (13, 1.25) to[american voltage source, l={$V_{ee}$}] (11.75, 1.25);
  	\draw (13, 1.25) -| (13.75, 2);
  	\node[ground, rotate=-90] at (11.75, 1.25){};
  	\draw (10.25, 10.75) -- (10.25, -0);
  	\draw (1, 9.25) -- (19, 9.25);
  	\node[shape=rectangle, minimum width=3.84cm, minimum height=1.09cm] at (4.813, 10.063){} node[anchor=north, align=center, text width=3.452cm, inner sep=6pt] at (4.813, 10.625){Full Schematic\\(Sink)};
  	\node[shape=rectangle, minimum width=3.84cm, minimum height=1.09cm] at (13.562, 10.063){} node[anchor=north, align=center, text width=3.452cm, inner sep=6pt] at (13.562, 10.625){Full Schematic\\(Source)};
  	\node[shape=rectangle, minimum width=3.09cm, minimum height=2.465cm] at (8.687, 7){} node[anchor=north west, align=left, text width=2.702cm, inner sep=6pt] at (7.125, 8.25){\textcolor{rgb,255:red,255;green,0;blue,0}{\footnotesize Exists Only For Analysis. In  Implementation take output directly from V\_\{out\}}};
  	\path[draw={rgb,255:red,255;green,0;blue,0}] (13.75, 2) -| (15, 2.25);
  	\node[shape=rectangle, minimum width=3.09cm, minimum height=2.465cm] at (17.708, 2.562){} node[anchor=north west, align=left, text width=2.702cm, inner sep=6pt] at (16.145, 3.812){\textcolor{rgb,255:red,255;green,0;blue,0}{\footnotesize Exists Only For Analysis. In  Implementation take output directly from V\_\{out\}}};
  	\node[shape=rectangle, minimum width=0.965cm, minimum height=0.465cm](N7) at (7, 3.71){} node[anchor=center] at (N7.text){$V_{x}$};
  	\draw (6, 3.71) |- (6.5, 3.71);
  	\node[shape=rectangle, minimum width=0.965cm, minimum height=0.465cm](N8) at (16, 6.25){} node[anchor=center] at (N8.text){$V_{x}$};
  	\draw (14.98, 6.25) |- (15.48, 6.25);
  	\node[shape=rectangle, minimum width=0.965cm, minimum height=0.465cm](N9) at (4.75, 3.96){} node[anchor=center] at (N9.text){$I_{G_1}$};
  	\node[shape=rectangle, minimum width=0.965cm, minimum height=0.465cm](N10) at (13.71, 6){} node[anchor=center] at (N10.text){$I_{G_1}$};
  	\draw (12.5, 4.562) -- (12.5, 3.75);
  	\draw (1, -0) -- (19, -0);
  	\draw (12.5, -7.063) to[american current source, l_={$g_{m_1}v_{gs}$}] (12.5, -8.563);
  	\draw (14, -7.063) to[european resistor, l_={$g_{o_1}$}] (14, -8.563);
  	\draw (12.5, -8.563) -- (14, -8.563);
  	\draw (14, -7.063) -- (12.5, -7.063);
  	\draw (12.5, -8.563) -- (11.25, -8.563);
  	\draw (11.25, -7.063) -- (10.5, -7.063);
  	\node[shape=rectangle, minimum width=0.465cm, minimum height=0.465cm] at (11.25, -7.313){} node[anchor=center, align=center, text width=0.077cm, inner sep=6pt] at (11.25, -7.313){\footnotesize G};
  	\node[shape=rectangle, minimum width=0.465cm, minimum height=0.465cm] at (11.25, -8.313){} node[anchor=center, align=center, text width=0.077cm, inner sep=6pt] at (11.25, -8.313){\footnotesize S};
  	\node[shape=rectangle, minimum width=0.465cm, minimum height=0.465cm] at (14.5, -7.313){} node[anchor=center, align=center, text width=0.077cm, inner sep=6pt] at (14.5, -7.313){\footnotesize D};
  	\draw (10.5, -7.063) -- (9.75, -7.063);
  	\draw (8.25, -7.063) to[american current source, l={$g_{m_0}v_{gs}$}] (8.25, -8.563);
  	\draw (6.75, -7.063) to[european resistor, l={$g_{o_0}$}] (6.75, -8.563);
  	\draw (8.25, -7.063) -- (6.75, -7.063);
  	\draw (8.25, -8.563) -- (6.75, -8.563);
  	\draw (6.75, -7.063) -| (6.75, -5.563);
  	\draw (9.75, -8.563) -- (8.25, -8.563);
  	\node[shape=rectangle, minimum width=0.465cm, minimum height=0.465cm] at (9.75, -7.313){} node[anchor=center, align=center, text width=0.077cm, inner sep=6pt] at (9.75, -7.313){\footnotesize G};
  	\node[shape=rectangle, minimum width=0.465cm, minimum height=0.465cm] at (9.75, -8.313){} node[anchor=center, align=center, text width=0.077cm, inner sep=6pt] at (9.75, -8.313){\footnotesize S};
  	\node[shape=rectangle, minimum width=0.465cm, minimum height=0.465cm] at (6.25, -7.313){} node[anchor=center, align=center, text width=0.077cm, inner sep=6pt] at (6.25, -7.313){\footnotesize D};
  	\draw (6.75, -1.25) to[european resistor, l={$G_4$}] (6.75, -2.75);
  	\draw (6.75, -2.75) -| (6.75, -3.25);
  	\draw (6.75, -1.25) -| (6.75, -0.75) -- (7.5, -0.75) -| (7.5, -1);
  	\draw (14, -2.063) to[european resistor, l={$G_5$}] (14, -3.563);
  	\draw (14, -3.563) -| (14, -4.063);
  	\draw (14, -2.063) -| (14, -1.563) |- (13.25, -1.563) -| (13.25, -1.813);
  	\draw (14, -4.063) -- (14.75, -4.063);
  	\node[shape=rectangle, minimum width=0.715cm, minimum height=0.465cm](N11) at (14.875, -6.313){} node[anchor=center] at (N11.text){$V_{x}$};
  	\node[shape=rectangle, minimum width=0.965cm, minimum height=0.465cm](N12) at (15.25, -4.063){} node[anchor=center] at (N12.text){$V_{out}$};
  	\node[ground] at (9.75, -8.563){};
  	\node[ground] at (11.25, -8.563){};
  	\node[ground] at (7.5, -1){};
  	\node[ground] at (13.25, -1.813){};
  	\node[shape=rectangle, minimum width=0.965cm, minimum height=0.465cm](N13) at (10, -6.375){} node[anchor=center] at (N13.text){$I_{G_1}$};
  	\node[shape=rectangle, minimum width=4.215cm, minimum height=1.09cm] at (10.25, -0.937){} node[anchor=north, align=center, text width=3.827cm, inner sep=6pt] at (10.25, -0.375){Small Signal Schematic\\(Sink)};
  	\draw (6, 3.52) -| (6, 3.71);
  	\node[nmos, arrowmos](N14) at (6, 4.48){} node[anchor=west] at (N14.text){$M_2$};
  	\draw (15, 6.398) |- (14.77, 6.25);
  	\node[shape=rectangle, minimum width=0.965cm, minimum height=0.465cm](N15) at (2.5, 3.71){} node[anchor=center] at (N15.text){$V_{y}$};
  	\draw (3, 3.71) |- (3.5, 3.71);
  	\node[pmos, arrowmos](N16) at (15, 5.48){} node[anchor=west] at (N16.text){$M_2$};
  	\node[shape=rectangle, minimum width=0.965cm, minimum height=0.465cm](N17) at (11.5, 6.25){} node[anchor=center] at (N17.text){$V_{y}$};
  	\draw (12, 6.25) |- (12.5, 6.25);
  	\draw (12.5, -4.063) to[american current source, l_={$g_{m_2}v_{gs}$}] (12.5, -5.563);
  	\draw (14, -4.063) to[european resistor, l_={$g_{o_2}$}] (14, -5.563);
  	\draw (12.5, -5.563) -- (14, -5.563);
  	\draw (14, -4.063) -- (12.5, -4.063);
  	\draw (12.5, -5.563) -- (11.25, -5.563);
  	\draw (11.25, -4.063) -- (10.5, -4.063);
  	\node[shape=rectangle, minimum width=0.465cm, minimum height=0.465cm] at (11.25, -4.313){} node[anchor=center, align=center, text width=0.077cm, inner sep=6pt] at (11.25, -4.313){\footnotesize G};
  	\node[shape=rectangle, minimum width=0.465cm, minimum height=0.465cm] at (11.25, -5.313){} node[anchor=center, align=center, text width=0.077cm, inner sep=6pt] at (11.25, -5.313){\footnotesize S};
  	\draw (14, -5.563) -- (14, -7.063);
  	\draw (10.5, -7.063) -| (10.5, -6.313) -- (14, -6.313);
  	\draw (14, -6.313) -- (14.5, -6.313);
  	\node[shape=rectangle, minimum width=0.715cm, minimum height=0.465cm](N18) at (5.875, -6.25){} node[anchor=center] at (N18.text){$V_{y}$};
  	\draw (6.25, -6.25) -- (6.75, -6.25);
  	\node[nmos, arrowmos, xscale=-1](N19) at (3.5, 4.48){} node[anchor=east] at (N19.text){$M_3$};
  	\draw (3.5, 3.71) -| (3.5, 3.52);
  	\draw (4.48, 4.48) -- (5.02, 4.48);
  	\draw (4.75, 4.5) |- (3.5, 5.5);
  	\draw (3.5, 5.25) -| (3.5, 5.5);
  	\draw (3.5, 6.04) -| (3.5, 5.5);
  	\node[shape=rectangle, minimum width=0.965cm, minimum height=0.465cm](N20) at (2.5, 5.5){} node[anchor=center] at (N20.text){$V_{z}$};
  	\draw (3, 5.5) |- (3.5, 5.5);
  	\node[shape=rectangle, minimum width=0.965cm, minimum height=0.465cm](N21) at (4.75, 5.75){} node[anchor=center] at (N21.text){$I_{G_2}$};
  	\node[pmos, arrowmos, xscale=-1](N22) at (12.5, 5.48){} node[anchor=east] at (N22.text){$M_3$};
  	\draw (13.48, 5.48) -- (14.02, 5.48);
  	\draw (12.5, 4.562) -- (12.5, 4.71);
  	\draw (12.5, 6.25) |- (12.5, 6.398);
  	\draw (13.75, 5.5) |- (12.5, 4.562);
  	\node[shape=rectangle, minimum width=0.965cm, minimum height=0.465cm](N23) at (13.75, 4.312){} node[anchor=center] at (N23.text){$I_{G_2}$};
  	\draw (10.5, -4.063) -- (9.75, -4.063);
  	\draw (8.25, -4.063) to[american current source, l={$g_{m_3}v_{gs}$}] (8.25, -5.563);
  	\draw (6.75, -4.063) to[european resistor, l={$g_{o_3}$}] (6.75, -5.563);
  	\draw (8.25, -4.063) -- (6.75, -4.063);
  	\draw (8.25, -5.563) -- (6.75, -5.563);
  	\draw (9.75, -5.563) -- (8.25, -5.563);
  	\node[shape=rectangle, minimum width=0.465cm, minimum height=0.465cm] at (9.75, -4.313){} node[anchor=center, align=center, text width=0.077cm, inner sep=6pt] at (9.75, -4.313){\footnotesize G};
  	\node[shape=rectangle, minimum width=0.465cm, minimum height=0.465cm] at (9.75, -5.313){} node[anchor=center, align=center, text width=0.077cm, inner sep=6pt] at (9.75, -5.313){\footnotesize S};
  	\node[shape=rectangle, minimum width=0.465cm, minimum height=0.465cm] at (6.25, -4.313){} node[anchor=center, align=center, text width=0.077cm, inner sep=6pt] at (6.25, -4.313){\footnotesize D};
  	\draw (10.5, -4.063) |- (6.75, -3.25) |- (6.75, -4.063);
  	\node[shape=rectangle, minimum width=0.965cm, minimum height=0.465cm](N24) at (11.5, 4.562){} node[anchor=center] at (N24.text){$V_{z}$};
  	\draw (12, 4.562) |- (12.5, 4.562);
  	\node[shape=rectangle, minimum width=0.715cm, minimum height=0.465cm](N25) at (5.875, -3.25){} node[anchor=center] at (N25.text){$V_{z}$};
  	\draw (6.25, -3.25) -- (6.75, -3.25);
  	\node[shape=rectangle, minimum width=0.965cm, minimum height=0.465cm](N26) at (10.5, -3){} node[anchor=center] at (N26.text){$I_{G_2}$};
  \end{tikzpicture}
\end{center}

\subsubsection{Transfer Function}
Applying KCL at $V_{out}$ on Full Schematic (for Sink):
\begin{align*}
  (V_{out} - V_{cc})G_5 + I_{M_2} &= 0 \\
  V_{out}G_5 - V_{cc}G_5 + I_{M_2} &= 0 \\
  V_{out} &= \frac{V_{cc}G_5 - I_{M_2}}{G_5} \tag{i}
\end{align*}
Applying KCL at $V_{x}$ on Full Schematic (for Sink):
\begin{align*}
  I_{M_1} + I_{G_1} - I_{M_2} &= 0 \\
  I_{M_1} + I_{G_1} &= I_{M_2} \\
  I_{M_1} &\approx I_{M_2} \quad \text{[assuming } I_{G_1} \approx 0 \text{ by MOS physics]} \tag{ii}
\end{align*}
Applying KCL at $V_{y}$ on Full Schematic (for Sink):
\begin{align*}
  I_{M_0} - I_{M_3} &= 0 \\
  I_{M_0} &= I_{M_3} \tag{iii}
\end{align*}
Applying KCL at $V_{z}$ on Full Schematic (for Sink):
\begin{align*}
  (V_{z} - V_{cc})G_4 + I_{M_3} + I_{G_2} &= 0 \\
  V_{z}G_4 - V_{cc}G_4 + (I_{M_3} + I_{G_2}) &= 0 \\
  V_{z}G_4 &= V_{cc}G_4 - (I_{M_3} + I_{G_2}) \\
  V_{z} &= V_{cc} - \frac{(I_{M_3} + I_{G_2})}{G_4} \tag{iv}
\end{align*}
Here's a table of the voltages at each terminal of the
transistors:\\
\begin{tabularx}{\linewidth}{|X|X|X|X|} \hline
  Transistor & Gate (G) & Source (S) & Drain (D) \\ \hline
  $M_0$ & $V_{x}$ & $V_{ee}$ & $V_{y}$ \\ \hline
  $M_1$ & $V_{x}$ & $V_{ee}$ & $V_{x}$ \\ \hline
  $M_2$ & $V_{z}$ & $V_{x}$ & $V_{out}$ \\ \hline
  $M_3$ & $V_{z}$ & $V_{y}$ & $V_{z}$ \\ \hline
\end{tabularx}\\
Here's a table of $V_{gs}$ and $V_{ds}$ for each transistor:\\
\begin{tabularx}{\linewidth}{|X|X|X|} \hline
  Transistor & $V_{gs}$ & $V_{ds}$ \\ \hline
  $M_0$ & $V_{x} - V_{ee}$ & $V_{y} - V_{ee}$ \\ \hline
  $M_1$ & $V_{x} - V_{ee}$ & $V_{x} - V_{ee}$ \\ \hline
  $M_2$ & $V_{z} - V_{x}$ & $V_{out} - V_{x}$ \\ \hline
  $M_3$ & $V_{z} - V_{y}$ & $V_{z} - V_{y}$ \\ \hline
\end{tabularx}\\
Using the above tables, we can express the currents through
each transistor as follows:
\begin{align*}
  I_{M_0} &= k_{n_0}(V_{gs_0} - V_{th})^2 \cdot (1 + \lambda V_{ds_0}) \\
          &= k_{n_0}((V_{x} - V_{ee}) - V_{th})^2 \cdot (1 + \lambda (V_{y} - V_{ee})) \\
          &= k_{n_0}(V_{x} - V_{ee} - V_{th})^2 \quad \text{[assuming } \lambda = 0 \text{]} \tag{v} \\ \\
  I_{M_1} &= k_{n_1}(V_{gs_1} - V_{th})^2 \cdot (1 + \lambda V_{ds_1}) \\
          &= k_{n_1}((V_{x} - V_{ee}) - V_{th})^2 \cdot (1 + \lambda (V_{x} - V_{ee})) \\
          &= k_{n_1}(V_{x} - V_{ee} - V_{th})^2 \quad \text{[assuming } \lambda = 0 \text{]} \tag{vi} \\ \\
  I_{M_2} &= k_{n_2}(V_{gs_2} - V_{th})^2 \cdot (1 + \lambda V_{ds_2}) \\
          &= k_{n_2}((V_{z} - V_{x}) - V_{th})^2 \cdot (1 + \lambda (V_{out} - V_{x})) \\
          &= k_{n_2}\left(\left(V_{cc} - \frac{(I_{M_3} + I_{G_2})}{G_4}\right) - V_{x} - V_{th}\right)^2 \cdot \left(1 + \lambda\left(\frac{V_{cc}G_5 - I_{M_2}}{G_5}\right) - V_{x}\right) \\
          &= k_{n_2}\left(\left(V_{cc} - \frac{(I_{M_3} + I_{G_2})}{G_4}\right) - V_{x} - V_{th}\right)^2 \quad \text{[assuming } \lambda = 0 \text{]} \tag{vii} \\ \\
  I_{M_3} &= k_{n_3}(V_{gs_3} - V_{th})^2 \cdot (1 + \lambda V_{ds_3}) \\
          &= k_{n_3}((V_{z} - V_{y}) - V_{th})^2 \cdot (1 + \lambda (V_{z} - V_{y})) \\
          &= k_{n_3}\left(\left(V_{cc} - \frac{(I_{M_3} + I_{G_2})}{G_4}\right) - V_{y} - V_{th}\right)^2 \codt \left(1 + \lambda\left(\left(V_{cc} - \frac{(I_{M_3} + I_{G_2})}{G_4}\right) - V_{y}\right)\right) \\
          &= k_{n_3}\left(\left(V_{cc} - \frac{(I_{M_3} + I_{G_2})}{G_4}\right) - V_{y} - V_{th}\right)^2 \quad \text{[assuming } \lambda = 0 \text{]} \tag{viii}
\end{align*}
It can be seen that,
\begin{align*}
  I_{in} &= I_{M_3} = I_{M_0} \quad \text{[from (iii)]} \\
  I_{out} &= I_{M_2} = I_{M_1} \quad \text{[from (ii)]} \\
  I_{M_0} &= I_{M_1} \quad \text{[from (v) and (vi) and assuming } k_{n_0} = k_{n_1} \text{]} \\
  \therefore  I_{out} &= I_{in}
\end{align*}
Thus,
\begin{align*}
  \boxed{\frac{I_{out}}{I_{in}} \approx \frac{I_{M_2}}{I_{M_3}} \approx \frac{I_{M_1}}{I_{M_0}} \approx \frac{k_{n_1}}{k_{n_0}} = \frac{\frac{\mu_n C_{ox} W_1}{L_1}}{\frac{\mu_n C_{ox} W_0}{L_0}} = \frac{\frac{W_1}{L_1}}{\frac{W_0}{L_0}} = \frac{W_1 L_0}{L_1 W_0}}
\end{align*}
The analysis is exactly the same for the Source
configuration, with the only difference being the
use of PMOS transistors instead of NMOS transistors.

\subsubsection{Analysis}
\begin{enumerate}
  \item \textbf{Setting Current to be Mirrored}:
    \begin{enumerate}
      \item Choose appropriate transistor dimensions
        ($W$ and $L$) for transistors $M_0$ and $M_1$
        to achieve the desired current gain.
        \begin{align*}
          \frac{I_{out}}{I_{in}} &= \frac{W_1 L_0}{L_1 W_0}
        \end{align*}
      \item Choose a suitable overdrive voltage value
        ($V_{ov}$) for computing the bias voltages at
        nodes $V_{x}$, $V_{y}$, and $V_{z}$.
        \begin{align*}
          V_{ov_0} &= V_{gs_0} - V_{th} \\
          V_{ov_0} &= V_{x} - V_{ee} - V_{th} \\
          V_{x} &= V_{ov_0} + V_{ee} + V_{th} \\
          \\
          V_{ov_2} &= V_{gs_2} - V_{th} \\
          V_{ov_2} &= V_{z} - V_{x} - V_{th} \\
          V_{z} &= V_{ov_2} + V_{x} + V_{th} \\
          \\
          V_{ov_3} &= V_{gs_3} - V_{th} \\
          V_{ov_3} &= V_{z} - V_{y} - V_{th} \\
          V_{y} &= V_{z} - V_{ov_3} - V_{th}
        \end{align*}
      \item Using the chosen $V_{ov}$, compute $W_0$
          to set the required bias current $I_{in}$.
          \begin{align*}
            I_{in} = I_{M_0} &= \frac{1}{2} \mu_n C_{ox} \frac{W_0}{L_0} V_{ov_0}^2 \\
            \therefore W_0 &= \frac{2 I_{in} L_0}{\mu_n C_{ox} V_{ov_0}^2}
          \end{align*}
          $W_1$ can then be computed using the current
          gain relation from step (a).
      \item Compute the required value of $G_4$ (or
          $R_4$) to ensure that the voltage at node
          $V_{z}$ is as calculated in step (b).
          \begin{align*}
            V_{z} &= V_{cc} - \frac{(I_{M_3} + I_{G_2})}{G_4} \\
            V_{z} &\approx V_{cc} - \frac{I_{M_3}}{G_4} \quad \text{[assuming } I_{G_2} \approx 0 \text{ by MOS physics]} \\
            V_{z} - V_{cc} &\approx - \frac{I_{M_3}}{G_4} \\
            \frac{(I_{M_3})}{(V_{cc} - V_{z})} &\approx G_4 \\
            \frac{(I_{in})}{(V_{cc} - V_{z})} &\approx G_4 \quad \text{[} I_{in} = I_{M_3} = I_{M_0} \text{]} \\
            \therefore R_4 &\approx \frac{V_{cc} - V_{z}}{I_{in}}
          \end{align*}
    \end{enumerate}
  \item \textbf{Output Load Limit}: Maximum possible
    load that can be connected at the output before
    the transistor $M_2$ goes out of saturation can
    be calculated as follows:
    \begin{align*}
      R_{out, max} &= \frac{V_{L, max}}{I_{M_2}} \quad \text{[where, } V_{ee} \leq V_{L, max} \leq V_{cc} \text{]}
    \end{align*}
  \item \textbf{Impedances}:
    \paragraph{At Input:} \textbf{Ignoring $G_4$
      which is used to set the current,} the input
      impedance is simply the output resistance of
      transistor $M_0$ scaled by the gain of
      transistor $M_3$.
      \begin{align*}
        Z_{in} &\approx \frac{1}{g_{o_0}\frac{g_{o_3}}{g_{m_3}}} \\
        &\approx \frac{1}{(\lambda_{M_0} I_{M_0}) \cdot \frac{(\lambda_{M_3} I_{M_3})}{\left(\frac{2 I_{M_3}}{V_{ov_3}}\right)}} \\
        &\approx \frac{1}{\lambda_{M_0} I_{M_0}} \cdot \frac{1}{\frac{\lambda_{M_3} V_{ov_3}}{2}} \\
        &\approx \frac{1}{\lambda_{M_0} I_{M_0}} \cdot \frac{2}{\lambda_{M_3} V_{ov_3}}\\
        &\approx \frac{2}{\lambda_{M_0} \lambda_{M_3} I_{M_0} (V_{gs_3} - V_{th})} \\
        &\approx \frac{2}{\lambda_{M_0} \lambda_{M_3} I_{M_0} ((V_{z} - V_{y}) - V_{th})}
      \end{align*}
    \paragraph{At Output:} \textbf{Ignoring $G_5$
      which is used as a load,} the output impedance
      is the output resistance of transistor $M_1$
      scaled by the gain of transistor $M_2$.
      \begin{align*}
        Z_{out} &\approx \frac{1}{g_{o_1}\frac{g_{o_2}}{g_{m_2}}} \\
        &\approx \frac{1}{(\lambda_{M_1} I_{M_1}) \cdot \frac{(\lambda_{M_2} I_{M_2})}{\left(\frac{2 I_{M_2}}{V_{ov_2}}\right)}} \\
        &\approx \frac{1}{\lambda_{M_1} I_{M_1}} \cdot \frac{1}{\frac{\lambda_{M_2} V_{ov_2}}{2}} \\
        &\approx \frac{1}{\lambda_{M_1} I_{M_1}} \cdot \frac{2}{\lambda_{M_2} V_{ov_2}}\\
        &\approx \frac{2}{\lambda_{M_1} \lambda_{M_2} I_{M_1} (V_{gs_2} - V_{th})} \\
        &\approx \frac{2}{\lambda_{M_1} \lambda_{M_2} I_{M_1} ((V_{z} - V_{x}) - V_{th})}
      \end{align*}
\end{enumerate}

\subsubsection{Example Design}

\paragraph{Question:}
Design a Modified Wilson Current Mirror in
Sink configuration to achieve a current gain of
2, with an input bias current of $I_{in} =
10\mu A$. Assume the following parameters for
the NMOS transistors used:
\begin{align*}
  \mu_n C_{ox} = 100 \mu A/V^2 \spc V_{th} = 0.7 V \spc \lambda = 0.02 V^{-1} \\
  V_{cc} = 5 V \scp V_{ee} = 0 V \spc L_0 = L_1 = L_2 = L_3 = 1 \mu m
\end{align*}
Find the output impedance and the maximum
possible load that can be connected at the output
assuming $V_{L,max} = V_{cc}$
\paragraph{Solution:}
\begin{enumerate}
  \item \textbf{Choose appropriate transistor dimensions:}
    \begin{align*}
      \frac{I_{out}}{I_{in}} &= \frac{W_1 L_0}{L_1 W_0} \\
      2 &= \frac{W_1 \cdot 1 \mu m}{1 \mu m \cdot W_0} \\
      W_1 &= 2 \cdot W_0
    \end{align*}
  \item \textbf{Choose a suitable overdrive voltage and determine
    the resulting bias voltages:}
    \begin{align*}
      V_{ov_0} &= 0.2 V \\
      V_{x} &= V_{ov_0} + V_{ee} + V_{th} = 0.2 + 0 + 0.7 = 0.9 V \\
      \\
      V_{ov_2} &= 0.2 V \\
      V_{z} &= V_{ov_2} + V_{x} + V_{th} = 0.2 + 0.9 + 0.7 = 1.8 V \\
      \\
      V_{ov_3} &= 0.2 V \\
      V_{y} &= V_{z} - V_{ov_3} - V_{th} = 1.8 - 0.2 - 0.7 = 0.9 V
    \end{align*}
  \item \textbf{Compute $W_0$ to set the required bias current:} For $I_{in} = 10 \mu A$
    \begin{align*}
      W_0 &= \frac{2 I_{in} L_0}{\mu_n C_{ox} V_{ov_0}^2} \\
          &= \frac{2 \cdot 10 \mu A \cdot 1 \mu m}{100 \mu A \cdot (0.2)^2} \\
          &= 5 \mu m
    \end{align*}
    Using (i) and (ii):
    \begin{align*}
      W_1 &= 2 \cdot W_0 = 2 \cdot 5 \mu m = 10 \mu m
    \end{align*}
  \item \textbf{Compute the required value of $G_4$ (or $R_4$):}
    \begin{align*}
      R_4 &\approx \frac{V_{cc} - V_{z}}{I_{in}} \\
          &= \frac{5 V - 1.8 V}{10 \mu A} \\
          &= 320 k\Omega
    \end{align*}
  \item \textbf{Output Load Limit:}
    \begin{align*}
      R_{out, max} &= \frac{V_{L, max}}{I_{M_2}} = \frac{5 V}{10 \mu A} = 500 k\Omega
    \end{align*}
  \item \textbf{Impedances:}
    \paragraph{At Input:}
    \begin{align*}
      Z_{in} &\approx \frac{2}{\lambda_{M_0} \lambda_{M_3} I_{M_0} (V_{gs_3} - V_{th})} \\
             &= \frac{2}{0.02 \imes 0.02 \times 10 \mu A \times (1.8 - 0.9)} \\
             &= 5.56 G\Omega
    \end{align*}
    \paragraph{At Output:}
    \begin{align*}
      Z_{out} &\approx \frac{2}{\lambda_{M_1} \lambda_{M_2} I_{M_1} (V_{gs_2} - V_{th})} \\
              &= \frac{2}{0.02 \imes 0.02 \times 20 \mu A \times (1.8 - 0.9)} \\
              &= 2.78 G\Omega
    \end{align*}
\end{enumerate}

\subsection{Cascode Mirror}

\subsubsection{Overview}
Constructed from four transistors, It improves
output resistance and accuracy by reducing the
effects of channel length modulation.

\subsubsection{Schematic}
\begin{center}
  \begin{tikzpicture}
  	% Paths, nodes and wires:
  	\node[shape=rectangle, draw={rgb,255:red,255;green,0;blue,0}, line width=1pt, dash pattern={on 4pt off 2pt on 1pt off 2pt on 1pt off 2pt}, minimum width=1.715cm, minimum height=2.195cm] at (6.125, 7.135){};
  	\draw (3.23, 8.04) to[american voltage source, l_={$V_{cc}$}] (1.98, 8.04);
  	\node[nmos, arrowmos, xscale=-1](N1) at (3.5, 2.75){} node[anchor=east] at (N1.text){$M_0$};
  	\node[nmos, arrowmos](N2) at (6, 2.75){} node[anchor=west] at (N2.text){$M_1$};
  	\draw (3.5, 7.54) to[european resistor, l={$G_4$}] (3.5, 6.04);
  	\draw (6, 5.79) |- (6, 5.25);
  	\draw (5.02, 2.75) -- (4.48, 2.75);
  	\draw (3.5, 1.98) -| (3.5, 1.75) -| (6, 1.98);
  	\draw (4, 1) to[american voltage source, l={$V_{ee}$}] (2.75, 1);
  	\draw (4, 1) -| (4.75, 1.75);
  	\node[ground, rotate=-90] at (1.98, 8.04){};
  	\node[ground, rotate=-90] at (2.75, 1){};
  	\draw[dash pattern={on 0.4pt off 0.8pt}, -stealth] (6, 5.5) -- (6.5, 5.5);
  	\node[shape=rectangle, minimum width=0.965cm, minimum height=0.465cm](N3) at (7, 5.46){} node[anchor=center] at (N3.text){$V_{out}$};
  	\draw (6, 7.5) to[european resistor, l={\textcolor{rgb,255:red,255;green,0;blue,0}{$G_5$}}] (6, 6);
  	\draw (6, 5.79) -- (6, 6.02);
  	\draw (3.5, 7.54) |- (3.23, 8.04);
  	\draw (3.5, 8.04) -| (6, 7.5);
  	\draw (10.25, 10.75) -- (10.25, -0);
  	\node[shape=rectangle, minimum width=5.59cm, minimum height=1.09cm] at (5.438, 9.937){} node[anchor=north, align=center, text width=5.202cm, inner sep=6pt] at (5.438, 10.5){Large Signal Schematic\\(Sink)};
  	\node[shape=rectangle, minimum width=3.09cm, minimum height=2.465cm] at (8.687, 7){} node[anchor=north west, align=left, text width=2.702cm, inner sep=6pt] at (7.125, 8.25){\textcolor{rgb,255:red,255;green,0;blue,0}{\footnotesize Exists Only For Analysis. In  Implementation take output directly from V\_\{out\}}};
  	\node[shape=rectangle, minimum width=0.965cm, minimum height=0.465cm](N4) at (7, 3.71){} node[anchor=center] at (N4.text){$V_{y}$};
  	\draw[dash pattern={on 0.4pt off 0.8pt}, -stealth] (6, 3.71) |- (6.5, 3.71);
  	\node[shape=rectangle, minimum width=0.965cm, minimum height=0.465cm](N5) at (4.75, 4.25){} node[anchor=center] at (N5.text){$I_{G_1}$};
  	\draw (17.625, 2.75) to[american current source, l_={$g_{m_1}v_{gs}$}] (17.625, 1.25);
  	\draw (19.125, 2.75) to[european resistor, l_={$g_{o_1}$}] (19.125, 1.25);
  	\draw (17.625, 1.25) -- (19.125, 1.25);
  	\draw (19.125, 2.75) -- (17.625, 2.75);
  	\draw (17.625, 1.25) -- (16.375, 1.25);
  	\draw (16.375, 2.75) -- (15.625, 2.75);
  	\node[shape=rectangle, minimum width=0.465cm, minimum height=0.465cm] at (16.375, 2.5){} node[anchor=center, align=center, text width=0.077cm, inner sep=6pt] at (16.375, 2.5){\footnotesize G};
  	\node[shape=rectangle, minimum width=0.465cm, minimum height=0.465cm] at (16.375, 1.5){} node[anchor=center, align=center, text width=0.077cm, inner sep=6pt] at (16.375, 1.5){\footnotesize S};
  	\node[shape=rectangle, minimum width=0.465cm, minimum height=0.465cm] at (19.625, 2.5){} node[anchor=center, align=center, text width=0.077cm, inner sep=6pt] at (19.625, 2.5){\footnotesize D};
  	\draw (15.625, 2.75) -- (14.875, 2.75);
  	\draw (13.375, 2.75) to[american current source, l={$g_{m_0}v_{gs}$}] (13.375, 1.25);
  	\draw (11.875, 2.75) to[european resistor, l={$g_{o_0}$}] (11.875, 1.25);
  	\draw (13.375, 2.75) -- (11.875, 2.75);
  	\draw (13.375, 1.25) -- (11.875, 1.25);
  	\draw (11.875, 2.75) -| (11.875, 4.25);
  	\draw (14.875, 1.25) -- (13.375, 1.25);
  	\node[shape=rectangle, minimum width=0.465cm, minimum height=0.465cm] at (14.875, 2.5){} node[anchor=center, align=center, text width=0.077cm, inner sep=6pt] at (14.875, 2.5){\footnotesize G};
  	\node[shape=rectangle, minimum width=0.465cm, minimum height=0.465cm] at (14.875, 1.5){} node[anchor=center, align=center, text width=0.077cm, inner sep=6pt] at (14.875, 1.5){\footnotesize S};
  	\node[shape=rectangle, minimum width=0.465cm, minimum height=0.465cm] at (11.375, 2.5){} node[anchor=center, align=center, text width=0.077cm, inner sep=6pt] at (11.375, 2.5){\footnotesize D};
  	\draw (11.875, 8.562) to[european resistor, l={$G_4$}] (11.875, 7.062);
  	\draw (11.875, 7.062) -| (11.875, 6.562);
  	\draw (11.875, 8.562) -| (11.875, 9.062) -- (12.625, 9.062) -| (12.625, 8.812);
  	\draw (19.125, 7.75) to[european resistor, l={\textcolor{rgb,255:red,255;green,0;blue,0}{$G_5$}}] (19.125, 6.25);
  	\draw (19.125, 6.25) -| (19.125, 5.75);
  	\draw (19.125, 7.75) -| (19.125, 8.25) |- (18.375, 8.25) -| (18.375, 8);
  	\node[shape=rectangle, minimum width=0.715cm, minimum height=0.465cm](N6) at (20, 3.5){} node[anchor=center] at (N6.text){$V_{y}$};
  	\node[shape=rectangle, minimum width=0.965cm, minimum height=0.465cm](N7) at (20.125, 5.75){} node[anchor=center] at (N7.text){$V_{out}$};
  	\node[ground] at (14.875, 1.25){};
  	\node[ground] at (16.375, 1.25){};
  	\node[ground] at (12.625, 8.812){};
  	\node[ground] at (18.375, 8){};
  	\node[shape=rectangle, minimum width=0.965cm, minimum height=0.465cm](N8) at (16.625, 3.563){} node[anchor=center] at (N8.text){$I_{G_1}$};
  	\node[shape=rectangle, minimum width=4.84cm, minimum height=1.09cm] at (15.25, 9.937){} node[anchor=north, align=center, text width=4.452cm, inner sep=6pt] at (15.25, 10.5){Small Signal Schematic\\(Sink)};
  	\draw (6, 3.52) -| (6, 3.71);
  	\node[nmos, arrowmos](N9) at (6, 4.48){} node[anchor=west] at (N9.text){$M_2$};
  	\node[shape=rectangle, minimum width=0.965cm, minimum height=0.465cm](N10) at (2.5, 3.71){} node[anchor=center] at (N10.text){$V_{x}$};
  	\draw[dash pattern={on 0.4pt off 0.8pt}, -stealth] (3.5, 3.71) |- (3, 3.71);
  	\draw (17.625, 5.75) to[american current source, l_={$g_{m_2}v_{gs}$}] (17.625, 4.25);
  	\draw (19.125, 5.75) to[european resistor, l_={$g_{o_2}$}] (19.125, 4.25);
  	\draw (17.625, 4.25) -- (19.125, 4.25);
  	\draw (19.125, 5.75) -- (17.625, 5.75);
  	\draw (17.625, 4.25) -- (16.375, 4.25);
  	\draw (16.375, 5.75) -- (15.625, 5.75);
  	\node[shape=rectangle, minimum width=0.465cm, minimum height=0.465cm] at (16.375, 5.5){} node[anchor=center, align=center, text width=0.077cm, inner sep=6pt] at (16.375, 5.5){\footnotesize G};
  	\node[shape=rectangle, minimum width=0.465cm, minimum height=0.465cm] at (16.375, 4.5){} node[anchor=center, align=center, text width=0.077cm, inner sep=6pt] at (16.375, 4.5){\footnotesize S};
  	\draw (19.125, 4.25) -- (19.125, 2.75);
  	\draw[dash pattern={on 0.4pt off 0.8pt}, -stealth] (19.125, 3.5) -- (19.625, 3.5);
  	\node[shape=rectangle, minimum width=0.715cm, minimum height=0.465cm](N11) at (11, 3.563){} node[anchor=center] at (N11.text){$V_{x}$};
  	\draw[dash pattern={on 0.4pt off 0.8pt}, -stealth] (11.875, 3.563) -- (11.375, 3.563);
  	\node[nmos, arrowmos, xscale=-1](N12) at (3.5, 4.48){} node[anchor=east] at (N12.text){$M_3$};
  	\draw (3.5, 3.71) -| (3.5, 3.52);
  	\draw (4.48, 4.48) -- (5.02, 4.48);
  	\draw (4.75, 4.5) |- (3.5, 5.5);
  	\draw (3.5, 5.25) -| (3.5, 5.5);
  	\draw (3.5, 6.04) -| (3.5, 5.5);
  	\node[shape=rectangle, minimum width=0.965cm, minimum height=0.465cm](N13) at (2.5, 5.5){} node[anchor=center] at (N13.text){$V_{z}$};
  	\draw[dash pattern={on 0.4pt off 0.8pt}, -stealth] (3.5, 5.5) |- (3, 5.5);
  	\node[shape=rectangle, minimum width=0.965cm, minimum height=0.465cm](N14) at (4.75, 6){} node[anchor=center] at (N14.text){$I_{G_2}$};
  	\draw (15.625, 5.75) -- (14.875, 5.75);
  	\draw (13.375, 5.75) to[american current source, l={$g_{m_3}v_{gs}$}] (13.375, 4.25);
  	\draw (11.875, 5.75) to[european resistor, l={$g_{o_3}$}] (11.875, 4.25);
  	\draw (13.375, 5.75) -- (11.875, 5.75);
  	\draw (13.375, 4.25) -- (11.875, 4.25);
  	\draw (14.875, 4.25) -- (13.375, 4.25);
  	\node[shape=rectangle, minimum width=0.465cm, minimum height=0.465cm] at (14.875, 5.5){} node[anchor=center, align=center, text width=0.077cm, inner sep=6pt] at (14.875, 5.5){\footnotesize G};
  	\node[shape=rectangle, minimum width=0.465cm, minimum height=0.465cm] at (14.875, 4.5){} node[anchor=center, align=center, text width=0.077cm, inner sep=6pt] at (14.875, 4.5){\footnotesize S};
  	\node[shape=rectangle, minimum width=0.465cm, minimum height=0.465cm] at (11.375, 5.5){} node[anchor=center, align=center, text width=0.077cm, inner sep=6pt] at (11.375, 5.5){\footnotesize D};
  	\draw (15.625, 5.75) |- (11.875, 6.562) |- (11.875, 5.75);
  	\node[shape=rectangle, minimum width=0.715cm, minimum height=0.465cm](N15) at (11, 6.562){} node[anchor=center] at (N15.text){$V_{z}$};
  	\draw[dash pattern={on 0.4pt off 0.8pt}, -stealth] (11.875, 6.562) -- (11.375, 6.562);
  	\node[shape=rectangle, minimum width=0.965cm, minimum height=0.465cm](N16) at (16.625, 6.562){} node[anchor=center] at (N16.text){$I_{G_2}$};
  	\draw[dash pattern={on 0.4pt off 0.8pt}, -stealth] (19.125, 5.75) -- (19.625, 5.75);
  	\draw[dash pattern={on 0.4pt off 0.8pt}, -stealth] (4.75, 5.5) -- (4.75, 5.75);
  	\draw[dash pattern={on 0.4pt off 0.8pt}, -stealth] (4.75, 3.71) -- (4.75, 4);
  	\draw[dash pattern={on 0.4pt off 0.8pt}, -stealth] (15.625, 6.562) -- (16.125, 6.562);
  	\draw[dash pattern={on 0.4pt off 0.8pt}, -stealth] (15.625, 3.563) -- (16.125, 3.563);
  	\node[shape=rectangle, draw={rgb,255:red,255;green,0;blue,0}, line width=1pt, dash pattern={on 4pt off 2pt on 1pt off 2pt on 1pt off 2pt}, minimum width=1.215cm, minimum height=2.195cm] at (19.375, 7.385){};
  	\draw (3.5, 3.71) -| (4.75, 2.75);
  	\draw (15.625, 2.75) |- (11.875, 3.563);
  \end{tikzpicture}
\end{center}

\subsubsection{Transfer Function}
Here's a table of the voltages at each terminal of the
transistors:\\
\begin{tabularx}{\linewidth}{|X|X|X|X|X|X|} \hline
  Transistor & Gate (G) & Source (S) & Drain (D) & $V_{gs}$ & $V_{ds}$ \\ \hline
  $M_0$ & $V_{x}$ & $V_{ee}$ & $V_{x}$ & $V_{x} - V_{ee}$ & $V_{x} - V_{ee}$ \\ \hline
  $M_1$ & $V_{x}$ & $V_{ee}$ & $V_{y}$ & $V_{x} - V_{ee}$ & $V_{y} - V_{ee}$ \\ \hline
  $M_2$ & $V_{z}$ & $V_{y}$ & $V_{out}$ & $V_{z} - V_{y}$ & $V_{out} - V_{y}$ \\ \hline
  $M_3$ & $V_{z}$ & $V_{x}$ & $V_{z}$ & $V_{z} - V_{x}$ & $V_{z} - V_{x}$ \\ \hline
\end{tabularx}\\

\begin{enumerate}
  \item \textbf{Large Signal Analysis}:\\
    Applying KCL at $V_{out}$:
    \begin{align*}
      (V_{out} - V_{cc})G_5 + I_{M_2} &= 0 \\
      V_{out}G_5 - V_{cc}G_5 + I_{M_2} &= 0 \\
      V_{out} = V_{cc} - \frac{I_{M_2}}{G_5} \tag{1}
    \end{align*}
    Applying KCL at $V_{x}$:
    \begin{align*}
      I_{M_0} + I_{G_1} - I_{M_3} &= 0 \\
      I_{M_0} + I_{G_1} = I_{M_3} \\
      I_{M_0} &\approx I_{M_3} \quad \text{[assuming } I_{G_1} \approx 0 \text{ by MOS physics]} \tag{2}
    \end{align*}
    Applying KCL at $V_{y}$:
    \begin{align*}
      I_{M_1} - I_{M_2} &= 0 \\
      I_{M_1} = I_{M_2} \tag{3}
    \end{align*}
    Applying KCL at $V_{z}$:
    \begin{align*}
      (V_{z} - V_{cc})G_4 + I_{M_3} + I_{G_2} &= 0 \\
      V_{z}G_4 - V_{cc}G_4 + I_{M_3} + I_{G_2} &= 0 \\
      V_{z}G_4 &= V_{cc}G_4 - (I_{M_3} + I_{G_2}) \\
      V_{z} &= V_{cc} - \frac{(I_{M_3} + I_{G_2})}{G_4} \tag{4}
    \end{align*}
    Using the transistor equations for $I_{M_0}$, $I_{M_1}$, $I_{M_2}$, and $I_{M_3}$:
    \begin{align*}
      I_{M_0} &= k_{n_0} (V_{gs_0} - V_{th})^2 (1 + \lambda V_{ds_0}) \\
              &= k_{n_0} (V_{x} - V_{ee} - V_{th})^2 (1 + \lambda (V_{x} - V_{ee})) \\
              &= k_{n_0} (V_{x} - V_{ee} - V_{th})^2 \quad \text{[assuming } \lambda \approx 0 \text{]} \tag{5} \\
      I_{M_1} &= k_{n_1} (V_{gs_1} - V_{th})^2 (1 + \lambda V_{ds_1}) \\
              &= k_{n_1} (V_{x} - V_{ee} - V_{th})^2 (1 + \lambda (V_{y} - V_{ee})) \\
              &= k_{n_1} (V_{x} - V_{ee} - V_{th})^2 \quad \text{[assuming } \lambda \approx 0 \text{]} \tag{6} \\
      I_{M_2} &= k_{n_2} (V_{gs_2} - V_{th})^2 (1 + \lambda V_{ds_2}) \\
              &= k_{n_2} (V_{z} - V_{y} - V_{th})^2 (1 + \lambda (V_{out} - V_{y})) \\
              &= k_{n_2} (V_{z} - V_{y} - V_{th})^2 \quad \text{[assuming } \lambda \approx 0 \text{]} \tag{7} \\
      I_{M_3} &= k_{n_3} (V_{gs_3} - V_{th})^2 (1 + \lambda V_{ds_3}) \\
              &= k_{n_3} (V_{z} - V_{x} - V_{th})^2 (1 + \lambda (V_{z} - V_{x})) \\
              &= k_{n_3} (V_{z} - V_{x} - V_{th})^2 \quad \text{[assuming } \lambda \approx 0 \text{]} \tag{8}
    \end{align*}
    It can be seen that,
    \begin{align*}
      I_{in} = I_{M_0} &= I_{M_3} \quad \text{[from (2)]} \\
      I_{out} = I_{M_1} &= I_{M_2} \quad \text{[from (3)]} \\
      I_{M_0} &= I_{M_1} \quad \text{[from (5) and (6) and assuming } k_{n_0} = k_{n_1} \text{]} \\
      \therefore, I_{in} &= I_{out} \\
    \end{align*}
    Thus,
    \begin{align*}
      \boxed{\frac{I_{out}}{I_{in}} \approx \frac{I_{M_2}}{I_{M_3}} \approx \frac{I_{M_1}}{I_{M_0}}
                                                                    \approx \frac{k_{n_1}}{k_{n_0}} = \frac{\frac{\mu_n C_{ox} W_1}{L_1}}{\frac{\mu_n C_{ox} W_0}{L_0}} = \frac{\frac{W_1}{L_1}}{\frac{W_0}{L_0}} = \frac{W_1 L_0}{L_1 W_0}
                                                                  = \frac{W_2 L_3}{L_2 W_3}} \\
    \end{align*}
  \item \textbf{Small Signal Analysis}:\\
    Applying KCL at $V_{out}$:
    \begin{align*}
      V_{out}G_5 + (V_{out} - V_{y})g_{o_2} + (V_{z} - V_{y})g_{m_2} &= 0 \\
      V_{out}G_5 + V_{out}g_{o_2} - V_{y}g_{o_2} + V_{z}g_{m_2} - V_{y}g_{m_2} &= 0 \\
      V_{out}(G_5 + g_{o_2}) - V_{y}(g_{o_2} + g_{m_2}) + V_{z}g_{m_2} &= 0 \tag{9}
    \end{align*}
    Applying KCL at $V_{x}$:
    \begin{align*}
      V_{x}g_{o_0} + I_{G_1} + V_{x}g_{m_0} + (V_{x} - V_{z})g_{o_3} - (V_{z} - V_{x})g_{m_3} &= 0 \\
      V_{x}g_{o_0} + V_{x}g_{m_0} + (V_{x} - V_{z})g_{o_3} - (V_{z} - V_{x})g_{m_3} &= 0 \quad \text{[as } I_{G} \approx 0 \text{ by MOS physics]} \\
      V_{x}g_{o_0} + V_{x}g_{m_0} + V_{x}g_{o_3} - V_{z}g_{o_3} - V_{z}g_{m_3} + V_{x}g_{m_3} &= 0 \\
      V_{x}(g_{o_0} + g_{m_0} + g_{o_3} + g_{m_3}) - V_{z}(g_{o_3} + g_{m_3}) &= 0 \tag{10}
    \end{align*}
    Applying KCL at $V_{y}$:
    \begin{align*}
      V_{y}g_{o_1} + V_{x}g_{m_1} + (V_{y} - V_{out})g_{o_2} - (V_{z} - V_{y})g_{m_2} &= 0 \\
      V_{y}g_{o_1} + V_{x}g_{m_1} + V_{y}g_{o_2} - V_{out}g_{o_2} - V_{z}g_{m_2} + V_{y}g_{m_2} &= 0 \\
      V_{y}(g_{o_1} + g_{o_2} + g_{m_2}) - V_{out}g_{o_2} - V_{z}g_{m_2} + V_{x}g_{m_1} &= 0 \tag{11}
    \end{align*}
    Applying KCL at $V_{z}$:
    \begin{align*}
      V_{z}G_4 + (V_{z} - V_{x})g_{o_3} + (V_{z} - V_{x})g_{m_3} + I_{G_2} &= 0 \\
      V_{z}G_4 + (V_{z} - V_{x})g_{o_3} + (V_{z} - V_{x})g_{m_3} &= 0 \\quad \text{[as } I_{G} \approx 0 \text{ by MOS physics]} \\
      V_{z}G_4 + V_{z}g_{o_3} - V_{x}g_{o_3} + V_{z}g_{m_3} - V_{x}g_{m_3} &= 0 \\
      V_{z}(G_4 + g_{o_3} + g_{m_3}) - V_{x}(g_{o_3} + g_{m_3}) &= 0 \tag{12}
    \end{align*}
    Adding voltage source $V_s$ at $V_z$ that supplies
    a current $I_s$ into $V_z$, we get:
    \begin{align*}
      V_{z}(G_4 + g_{o_3} + g_{m_3}) - V_{x}(g_{o_3} + g_{m_3}) - I_s &= 0 \tag{13}
    \end{align*}
    Fifth equation arises from Voltage source:
    \begin{align*}
      V_{z} = V_s \tag{14}
    \end{align*}
    Equations (9), (10), (11), (13), and (14) can be
    solved analytically. It is cumbersome so, I've
    chosen to give the final result directly:
    \begin{align*}
      \frac{V_{out}}{V_{s}} \quad \text{where,} \\
      V_{out} = - g_{m_0}g_{m_2}g_{o_1} - g_{m_1}g_{m_2}(g_{m_3} + g_{o_3}) - g_{m_1}g_{o_2}(g_{m_3} + g_{o_3}) - g_{m_2}g_{o_1}(g_{m_3} + (g_{o_0} + g_{o_3})) \\
      V_{s} = G_{5}[(g_{m_0}(g_{m_2} + g_{o_1} + g_{o_2}) + g_{m_2}(g_{m_3} + g_{o_0} + g_{o_3}) + g_{m_3}(g_{o_1} + g_{o_2})) + g_{o_0}(g_{o_1} + g_{o_2}) \\
              + (g_{o_3}(g_{o_1} + g_{o_2}))] + g_{o_1}g_{o_2}(g_{m_0} + g_{m_3} + g_{o_0} + g_{o_3})
    \end{align*}
\end{enumerate}


\subsubsection{Analysis}
\begin{enumerate}
  \item \textbf{Output Load (Max)}: Maximum possible
    load that can be connected at the output before
    the transistor $M_2$ goes out of saturation can
    be calculated as follows:
    \begin{align*}
      R_{out, max} &= \frac{V_{L, max}}{I_{M_2}} \quad \text{[where, } V_{ee} \leq V_{L, max} \leq V_{cc} \text{]}
    \end{align*}
  \item \textbf{Impedances}:
    \paragraph{At Input:}Applying a test current $I_{test}$
      at the input node $V_{z}$, from equation (12):
      \begin{align*}
        V_{z}(G_4 + g_{o_3} + g_{m_3}) - V_{x}(g_{o_3} + g_{m_3}) - I_{test} &= 0 \\
        V_{z}(G_4 + g_{o_3} + g_{m_3}) - V_{x}(g_{o_3} + g_{m_3}) &= I_{test} \tag{15}
      \end{align*}
      Solving equations (9), (10), (11), (15) simultaneously for $\frac{V_{z}}{I_{test}}$
      gives the input impedance. It is cumbersome so, I've chosen to give the final
      result directly from Octave:
      \begin{align*}
        Z_{in} = \frac{V_{z}}{I_{test}} = \frac{g_{m_0} + g_{m_3} + g_{o_0} + g_{o_3}}{G_{4}(g_{m_3} + g_{o_3}) + (g_{m_0} + g_{o_0})(G_{4} + g_{m_3} + g_{o_3})}
      \end{align*}
    \paragraph{At Output:}Applying a test current $I_{test}$
      at the output node $V_{out}$, from equation (9):
      \begin{align*}
        V_{out}(G_5 + g_{o_2}) - V_{y}(g_{o_2} + g_{m_2}) + V_{z}g_{m_2} - I_{test} &= 0 \\
        V_{out}(G_5 + g_{o_2}) - V_{y}(g_{o_2} + g_{m_2}) + V_{z}g_{m_2} &= I_{test} \tag{16}
      \end{align*}
      Solving equations (10), (11), (12) and (16) simultaneously for $\frac{V_{out}}{I_{test}}$
      gives the output impedance. It is cumbersome so, I've chosen to give the final
      result directly from Octave:
      \begin{align*}
        Z_{out} = \frac{V_{out}}{I_{test}} = \frac{g_{m_2} + g_{o_1} + g_{o_2}}{G_{5}(g_{m_2} + g_{o_1} + g_{o_2}) + g_{o_1} g_{o_2}}
      \end{align*}
\end{enumerate}


\subsubsection{Design}
Large Signal Analysis can be used to set the bias
currents and voltages in this circuit.
\begin{enumerate}
  \item Choose appropriate transistor dimensions
    ($W$ and $L$) for transistors $M_0$ and $M_1$
    to achieve the desired current gain.
    \begin{align*}
      \frac{I_{out}}{I_{in}} &= \frac{W_1 L_0}{L_1 W_0} = \frac{W_2 L_3}{L_2 W_3} \\
    \end{align*}
  \item Choose a suitable overdrive voltage value
    ($V_{ov}$) for computing the bias voltages at
    nodes $V_{x}$, $V_{y}$, and $V_{z}$.
    \begin{align*}
      V_{ov_0} &= V_{gs_0} - V_{th} \\
      V_{ov_0} &= V_{x} - V_{ee} - V_{th} \\
      V_{x} &= V_{ov_0} + V_{ee} + V_{th} \\
      \\
      V_{ov_3} &= V_{gs_3} - V_{th} \\
      V_{ov_3} &= V_{z} - V_{x} - V_{th} \\
      V_{z} &= V_{ov_3} + V_{x} + V_{th} \\
      \\
      V_{ov_2} &= V_{gs_2} - V_{th} \\
      V_{ov_2} &= V_{z} - V_{y} - V_{th} \\
      V_{y} &= V_{z} - V_{ov_2} - V_{th}
    \end{align*}
  \item Using the chosen $V_{ov}$, compute $W_0$
    to set the required bias current $I_{in}$.
    \begin{align*}
      I_{in} &\approx I_{M_0} = \frac{1}{2} \mu_n C_{ox} \frac{W_0}{L_0} V_{ov_0}^2 \\
             &\approx I_{M_3} = \frac{1}{2} \mu_n C_{ox} \frac{W_3}{L_3} V_{ov_3}^2 \\
      \therefore W_0 &= \frac{2 I_{in} L_0}{\mu_n C_{ox} V_{ov_0}^2} \\
                 W_3 &= \frac{2 I_{in} L_3}{\mu_n C_{ox} V_{ov_3}^2}
    \end{align*}
    $W_1$ can then be computed using the current
    gain relation from step (a). Similarly, $W_2$
    can be computed using the current gain relation
    from step (a).
  \item Compute the required value of $G_4$ (or
    $R_4$) to ensure that the voltage at node
    $V_{z}$ is as calculated in step (b).
    \begin{align*}
      V_{z} &= V_{cc} - \frac{(I_{M_3} + I_{G_2})}{G_4} \\
      V_{z} &\approx V_{cc} - \frac{I_{M_3}}{G_4} \quad \text{[assuming } I_{G_2} \approx 0 \text{ by MOS physics]} \\
      V_{z} - V_{cc} &\approx - \frac{I_{M_3}}{G_4} \\
      \frac{(I_{M_3})}{(V_{cc} - V_{z})} &\approx G_4 \\
      \frac{(I_{in})}{(V_{cc} - V_{z})} &\approx G_4 \quad \text{[} I_{in} = I_{M_3} = I_{M_0} \text{]} \\
      \therefore R_4 &\approx \frac{V_{cc} - V_{z}}{I_{in}}
    \end{align*}
\end{enumerate}

\subsubsection{Question}
Design a cascode current mirror with a current
gain of 5. The input current is $I_{in} = 10 \mu$A.
Assume the following specificiatons:
\begin{enumerate}
  \item $V_{cc} = +5$ V, $V_{ee} = 0$ V
  \item $V_{th} = 0.7$ V
  \item $\mu_n C_{ox} = 100 \mu A/V^2$
  \item $\lambda = 0.02$ V$^{-1}$
  \item $L_{min} = 1 \mu m$
  \item $V_{L, max} = V_{cc}$
  \item $V_{ov} = 0.2$ V for all transistors
\end{enumerate}
Find the following:
\begin{enumerate}
  \item Transistor dimensions ($W$ and $L$) for all
    transistors.
  \item Values of resistor $R_4$ required to set
    the bias current.
  \item Output voltage range and maximum output
    load that can be connected at the output.
  \item Small signal output resistance of the
    current mirror. Assume a load resistance
    $R_L = 100 k\Omega$ is connected at the
    output.
\end{enumerate}
\paragraph{Solution:}
\begin{enumerate}
  \item Transistor dimensions:
    \begin{align*}
      \frac{I_{out}}{I_{in}} &= \frac{W_1 L_0}{L_1 W_0} = \frac{W_2 L_3}{L_2 W_3} = 5 \\
      \\
      V_{x} &= V_{ov_0} + V_{ee} + V_{th} = 0.2 + 0 + 0.7 = 0.9 \text{ V} \\
      V_{z} &= V_{ov_3} + V_{x} + V_{th} = 0.2 + 0.9 + 0.7 = 1.8 \text{ V} \\
      V_{y} &= V_{z} - V_{ov_2} - V_{th} = 1.8 - 0.2 - 0.7 = 0.9 \text{ V} \\
      \\
      W_0 &= \frac{2 I_{in} L_0}{\mu_n C_{ox} V_{ov_0}^2} = \frac{2 \times 10 \times 10^{-6} \times 1 \times 10^{-6}}{100 \times 10^{-6} \times (0.2)^2} = 5 \mu m \\
      W_3 &= \frac{2 I_{in} L_3}{\mu_n C_{ox} V_{ov_3}^2} = \frac{2 \times 10 \times 10^{-6} \times 1 \times 10^{-6}}{100 \times 10^{-6} \times (0.2)^2} = 5 \mu m \\
      W_1 &= 5 \times \frac{W_0}{L_0} \times L_1 = 5 \times \frac{5}{1} \times 1 = 25 \mu m \\
      W_2 &= 5 \times \frac{W_3}{L_3} \times L_2 = 5 \times \frac{5}{1} \times 1 = 25 \mu m
    \end{align*}
  \item Value of resistor $R_4$:
    \begin{align*}
      R_4 &\approx \frac{V_{cc} - V_{z}}
                      {I_{in}} = \frac{5 - 1.8}{10 \times 10^{-6}} = 320 k\Omega
    \end{align*}
  \item Output voltage range and maximum output load:
    \begin{align*}
      V_{out, min} &\approx V_{y} + V_{ov_2} = 0.9 + 0.2 = 1.1 \text{ V} \\
      V_{out, max} &\approx V_{cc} = 5 \text{ V} \\
      R_{out, max} &= \frac{V_{L, max}}{I_{M_2}} = \frac{5 - 0}{10 \times 10^{-6}} = 500 k\Omega
    \end{align*}
  \item Small signal output resistance:
    \begin{align*}
      g_{m_2} &= \frac{2 I_{M_2}}{V_{ov_2}} = \frac{2 \times 10 \times 10^{-6}}{0.2} = 0.0001 \text{ S} \\
      g_{o_1} &= I_{M_1} \lambda = 10 \times 10^{-6} \times 0.02 = 0.2 \times 10^{-6} \text{ S} \\
      g_{o_2} &= I_{M_2} \lambda = 10 \times 10^{-6} \times 0.02 = 0.2 \times 10^{-6} \text{ S} \\
      G_5 &= \frac{1}{R_L} = \frac{1}{100 k\Omega} = 10^{-5} \text{ S} \\
      \\
      Z_{out} &= \frac{g_{m_2} + g_{o_1} + g_{o_2}}{G_{5}(g_{m_2} + g_{o_1} + g_{o_2}) + g_{o_1} g_{o_2}} \\
              &= \frac{0.0001 + 0.2 \times 10^{-6} + 0.2 \times 10^{-6}}{10^{-5}(0.0001 + 0.2 \times 10^{-6} + 0.2 \times 10^{-6}) + (0.2 \times 10^{-6})(0.2 \times 10^{-6})} \\
              &\approx 9.99 M\Omega
    \end{align*}
\end{enumerate}

\subsubsection{Code}
\begin{lstlisting}[language=Octave]
  clear; clc;
  pkg load symbolic;

  syms Vx Vy Vz Vout Vs Is Itest
  syms G5 G4 go0 gm0 go1 gm1 go2 gm2 go3 gm3

  % system kcl
  eq1 = Vout*(G5 + go2) - Vy*(go2 + gm2) + Vz*gm2 == 0; % KCL at Vout
  eq2 = Vx*(go0 + gm0 + go3 + gm3) - Vz*(go3 + gm3) == 0; % KCL at Vx
  eq3 = Vy*(go1 + go2 + gm2) - Vout*go2 - Vz*gm2 + Vx*gm1 == 0; % KCL at Vy
  eq4 = Vz*(G4 + go3 + gm3) - Vx*(go3 + gm3) == 0; % KCL at Vz

  % voltage transfer function: Vs forced at Vz
  eq4_vt = Vz*(G4 + go3 + gm3) - Vx*(go3 + gm3) - Is == 0;
  eq5_vt = Vz == Vs;
  sol_vt = solve([eq1, eq2, eq3, eq4_vt, eq5_vt], [Vout, Vx, Vy, Vz, Vs]);
  volt_transfer = simplify(sol_vt.Vout/sol_vt.Vs);
  % volt_transfer = collect(volt_transfer, G5);
  % volt_transfer = subs(volt_transfer, 'G4', 0);
  % volt_transfer = subs(volt_transfer, 'G5', 0);
  printf('=== Voltage Transfer Results ===\n');
  pretty(volt_transfer)
  printf('\n');

  % current transfer function: Is injected into Vz
  sol_ct = solve([eq1, eq2, eq3, eq4_vt], [Vout, Vx, Vy, Vz]);
  Iout = sol_ct.Vout * G5;
  curr_transfer = simplify(Iout/Is);
  % curr_transfer = collect(curr_transfer, G4*G5);
  printf('=== Current Transfer Results ===\n');
  pretty(curr_transfer)
  printf('\n');

  % input impedence: Itest injected into Vz
  eq4_zin = Vz*(G4 + go3 + gm3) - Vx*(go3 + gm3) - Itest == 0;
  sol_zin = solve([eq1, eq2, eq3, eq4_zin], [Vout, Vx, Vy, Vz]);
  zin = simplify(sol_zin.Vz/Itest);
  % zin = collect(zin, G4*G5);
  % zin = collect(zin, G4);
  % zin = subs(zin, 'G5', 0);
  printf('=== Input Impedence Results ===\n');
  pretty(zin);
  printf('\n');

  % output impedence: Itest injected into Vout
  eq1_zout = Vout*(G5 + go2) - Vy*(go2 + gm2) + Vz*gm2 - Itest == 0;
  sol_zout = solve([eq1_zout, eq2, eq3, eq4], [Vout, Vx, Vy, Vz]);
  zout = simplify(sol_zout.Vout/Itest);
  % zout = collect(zout, G4*G5);
  % zout = subs(zout, 'G4', 0);
  % zout = collect(zout, G5);
  printf('=== Output Impedence Results ===\n');
  pretty(zout)
  printf('\n');
\end{lstlisting}

